\documentclass[11pt,reqno]{amsart}
\usepackage{amsmath,amssymb,amsthm}
\usepackage[margin=1.1in]{geometry}
\usepackage{hyperref}
\usepackage{xr}
\externaldocument[M-]{erdos-634}
\externaldocument[C-]{erdos-634-companion}
\hypersetup{colorlinks=true,linkcolor=blue,citecolor=blue,urlcolor=blue}
\usepackage{booktabs}
\theoremstyle{plain}
\newtheorem{theorem}{Theorem}
\newtheorem{lemma}[theorem]{Lemma}
\newtheorem{corollary}[theorem]{Corollary}
\newtheorem{proposition}[theorem]{Proposition}
\theoremstyle{definition}
\newtheorem{remark}[theorem]{Remark}
\newtheorem{conjecture}[theorem]{Conjecture}
\newtheorem{hypothesis}[theorem]{Hypothesis}
\newtheorem{definition}[theorem]{Definition}
\newcommand{\Z}{\mathbb{Z}}
\newcommand{\Q}{\mathbb{Q}}
\newcommand{\R}{\mathbb{R}}
\newcommand{\N}{\mathbb{N}}

\title{Erd\H{o}s Problem 634: obstructions, scope limits, and closed routes}
\author{Vico Bonfioli}
\date{\today}


%% Rule 0 label macro: exactly one label per numbered statement.
\newcommand{\lab}[1]{\hfill\textnormal{\footnotesize\textsc{[#1]}}}

\begin{document}
\begin{abstract}
This note collects the negative results of the base-$\beta$ programme: the exact scope of each theorem
in the companion, the mechanisms proved unable to close the remaining case, and the routes closed with
their precise failure points. Nothing here is needed to read \cite{Main} or \cite{Companion}. It exists
so that closed routes are not retried, and so that each positive result carries its limits somewhere
citable rather than inline.
\end{abstract}
\maketitle

\section{How to read this}

Each entry states a limit or a closure. Labels mirror the companion: a remark labelled \texttt{O-foo}
here is the discussion belonging to \texttt{rem:foo} there, where a one-sentence summary is retained in
place so that no theorem stands without its scope. Numbering is independent of both other papers.

\section{Scope limits of the positive results}

\begin{remark}[Scope and consequences]\label{O-ray}

Thickness is used exactly once, and sharply: the next solution would be $n_a=2f-e$, which needs
$f\ge2e+1$ --- the thin regime. So in the regime that is open, the entire mismatch ray is now pinned
on \emph{both} sides, with no freedom beyond the order of the near-side edges. Explicitly:
$N=23$ has far $5^3$, near $\{6,9\}$; $N=59$ far $9^5$, near $\{20,25\}$; $N=66$ far $16^5$, near
$\{15^2,25^2\}$; $N=83$ far $11^6$, near $\{30,36\}$; $N=167$ far $39^8$, near $\{40^3,64^3\}$.
The count of interior junctions is therefore also determined --- $f-1$ on the far side and
$2(f-e)-1$ on the near side, none shared (Theorem~\ref{C-thm:align}(i)) --- so a hypothetical thick
tiling carries exactly $3f-2e-2$ T-junctions along this ray, each a $\pi$-vertex. The arithmetic is
machine-checked (\texttt{near\_side\_unique}, \texttt{BaseBetaWalks.lean}).
\end{remark}

\begin{remark}[Scope, and why $m=1$ is the rigid case]\label{O-walkscope}

Theorem~\ref{C-thm:walks}(i) generalizes the $e=1$ base analysis used for
Theorem~\ref{C-thm:basebeta-e1} to the infinite family $f>2e$ --- $e=1$, $f\ge3$; $e=2$, $f\ge5$;
$e=4$, $f\ge9$; \dots --- covering the primes $N=47,71,107,191,227,239,359,\dots$, and its parts (iii)--(iv),
which need no thinness, deliver the $e\ge2$ analogue of the $e=1$ structure theorem of
Remark~\ref{M-rem:lem6} for \emph{every} $(e,f)$: the equal sides are $b$-free and the base's $b$-count
is bounded by $j(f-e)\le e-1$. The step that fails for $m\ge2$ is the very first one: there
$Q\equiv em\pmod f$, so $j$ may be \emph{negative} and the level rises to $2em$. This is not an
artifact. The genuine $44$-tiling ($e{=}1,f{=}2,m{=}2$) has base walk $aaaaccca$, i.e.
$(P,Q,R)=(5,0,3)$, which is the parameter $j=-1$, $p=3$; and the $99$-tiling
($e{=}1,f{=}2,m{=}3$) has base walk $aabbbbbbbcc$, i.e. $(2,7,2)$, the parameter $j=2$, $p=0$. Both
satisfy $R\ge1$, as the $\gamma$-trap requires.
\end{remark}

\begin{remark}[The blocking hypothesis cannot be weakened to ``both ends are vertices'']\label{O-noedgetoedge}

It is worth recording what the step left open in Remark~\ref{C-rem:route1uniform} may \emph{not} appeal
to. In Lemma~\ref{C-lem:cchord} the hypothesis that the chord is blocked at both ends is what converts
it into a union of whole tile edges; one cannot replace it by the weaker hypothesis that both ends
are vertices of the tiling, because the tilings in this family are not edge-to-edge. The certified
$99$-tiling of member $(1,2)$ settles this concretely. Its $99$ tiles span $216$ distinct edges, and
$52$ of those carry a vertex of the tiling in their relative interior --- each exactly one. Moreover
its $89$ vertices contain $19$ pairs at distance exactly $a$, $23$ at distance exactly $b$ and $45$
at distance exactly $c$ that span no tile edge at all. (These counts are read off the coordinates of
the machine-checked certificate; in that embedding every $x$ is rational and every $y$ a rational
multiple of $\sqrt{15}$, so the chart $(x,y/\sqrt{15})$ is an affine rescaling of the plane and
decides collinearity and betweenness faithfully.)

Hence knowing that a segment has both endpoints at tiling vertices and has the length of a tile side
carries \emph{no} information about how it meets the tiling. Since $V$ and $E$ are vertices at
distance exactly $a$, this is precisely the shape of the hypothesis one is tempted to use, and it is
not available. Whatever forces $[V,E]$ to be a whole $a$-edge must consume the cascade's tile data
--- which tiles have been forced and where --- and not merely its vertex set. The same remark
explains why the tile-interior blocking of Remark~\ref{C-rem:tib}, and not a vertex-incidence
argument, is the mechanism the cascade actually runs on.
\end{remark}

\begin{remark}[A gap in the induction of Theorem~\ref{C-thm:n1}]\label{O-n1gap}

The induction step of Theorem~\ref{C-thm:n1} argues that at $V_k=(kc,0)+a\,u$ the tiles $T_k$ and
$P_{k-1}$ present $\gamma$ and $\beta$, ``leaving $\alpha$''. That is the straight-junction identity
$\pi-\gamma-\beta=\alpha$, which follows from $3\alpha+2\beta=\pi$ and $\gamma=2\alpha+\beta$. It is
legitimate at $k=0$, where $V_0=a\,u$ lies on an equal side and the junction is genuinely straight.
For $k\ge1$ the point $V_k$ is \emph{strictly interior}: the angle around it is $2\pi$, and
\[
2\pi-\gamma-\beta \;=\; 4\alpha+2\beta \;=\; \pi+\alpha \;\ne\;\alpha .
\]
The discrepancy is exactly $\pi$. For instance at $(e,f)=(2,5)$ one has
$V_1=(30.68,8.23)$, interior to the target $(0,0),(142,0),f^3u$; the residue there is $203.07^\circ$,
not the $\alpha=23.07^\circ$ the step requires. The same holds at every member we computed,
$(2,3)$, $(3,8)$, $(1,4)$ and $(3,7)$ among them. With the residue $\pi+\alpha$ the partner $P_k$ is
not forced, and the chain stops after $T_0,P_0,T_1$; that much is exactly
Proposition~\ref{M-prop:cornerpara}.

At $e=1$ the theorem is nevertheless true, by a different and much shorter route: the column
$(0,e,2e)=(0,1,2)$ has a base word of only three letters, and
Proposition~\ref{M-prop:cornerpara} places the first two and the last two base edges in $\{a,c\}$,
which for a three-letter word is all of them; the word then carries no $b$-edge, contradicting
$y=1$. That argument uses no induction and no separation.

At $e\ge2$ the word has $3e\ge6$ letters, the four pinned positions leave $3e-4\ge e$ interior slots,
and the $b$-edges fit inside them: at $(2,5)$ the word $c\,c\,b\,b\,c\,c$ satisfies
$25+25+21+21+25+25=142=e(3f^2-e^2)$ and survives every filter we have. So Theorem~\ref{C-thm:n1} is
established at $e=1$ only, and its conclusion at $e\ge2$ --- on which
Corollary~\ref{C-cor:basewalls} depends --- must be regarded as \textbf{open}. We have not refuted the
conclusion; we refute only the proof. The missing ingredient is sharp: force the $\beta$-tile at
$V_k$, whose flanks are $\{a,c\}$, to lay its $c$ rather than its $a$ along $[V_k,V_{k+1}]$. This is
an instance of the crossing question of Remark~\ref{C-rem:straddle}, which
Remark~\ref{rem:chainobstruction} already identifies as the obstruction, and
Remark~\ref{C-rem:whyf3} already warns that a pinning argument must say against what the ends are
pinned. Here the would-be blocking edge at the upper end is $P_k$'s own $c$-edge, the very thing
being derived.
\end{remark}

\begin{remark}[The $V_1$ step is repairable, and its arithmetic is machine-checked]\label{O-v1repair}\lab{CONJECTURE: complete proof sketch; the arithmetic cores VERIFIED}
The induction step \emph{is} recoverable at $k=1$, where --- and only where --- both relevant
anchors are boundary points: the west end $V_0=a\,u$ of the level segment and the far tip
$W=2a\,u$ of the chord's extension both lie on the equal side, so every escape of the interior
analysis, each of which must overrun one of them, is blocked. The closures then die on three
arithmetic facts, formalized and axiom-clean (\texttt{lean/Erdos634/V1Gaps.lean}): $f^2-2e^2$ is
not representable in $\langle a,b,c\rangle$ for coprime $f>2e$; no whole-edge prefix $\sigma$
and covering edge $L$ satisfy $\sigma+L=2b-xa$ with the overrun strictly inside $L$; and at
$e\ge2$ the length $c$ has the unique representation $(0,0,1)$ --- precisely the statement that
fails at $e=1$, where $c=fa$ is the crossing configuration itself. Consequently, at proof-sketch
rigor: the column $(0,e,2e)$ is excluded at $e=2$ for every coprime $f>4$. The count is sharper
than a mirror argument needs. At $e=2$ the word has six letters, two $b$ and four $c$, and no
$a$; Proposition~\ref{M-prop:cornerpara}, which is two-sided and needs no separation, pins the
first two and last two positions to $\{a,c\}$, so with no $a$ available those four positions
consume all four $c$s and the word is $ccbbcc$. The repair forces the third position to be $c$ as
well --- a fifth $c$ where four exist. The reflection argument is therefore \emph{not} needed at
$e=2$ (it is needed at $e=3$, where the same forcing leaves $c^3b^3c^3$ consistent). Two scope
notes: no new range in $f$ is gained, since $f>2e$ and separation both give $f\ge5$ at $e=2$ ---
what is gained is that the dichotomy step there no longer cites Hypothesis~\ref{C-hyp:walls};
and \emph{formalized} describes the number theory alone, every geometric step of the repair
lying on the far side of the obligations named in \cite{Companion}'s
Remark~\ref{C-rem:pingaps}. At $e=3$ the same forcing pins the column word to
$c^3b^3c^3$, and the $V_2$ step is the named residue. That step cannot be reached by refining the
present argument, and it is worth saying why. The chord tip at step $k$ is
$W_k=((k-1)c,0)+2a\,u$, which lies on the equal side's ray exactly when $k=1$: the offset
$((k-1)c,0)$ is horizontal while $u$ is not, so it can be absorbed only when it vanishes
(\texttt{V1Assembly.chord\_tip\_on\_ray\_iff}; numerically, $W_1$ is at distance $0$ from the
side at every member checked and $W_2$ is interior by roughly a tile diameter,
\texttt{code/analysis/chord\_tip\_check.py}). Any blocker at $V_2$ must therefore be a tile
already forced rather than the target's edge --- which is the locating problem in its
per-configuration form. A search at $(3,7)$ on the surviving word $c^3b^3c^3$ reached five million
nodes without exhausting, an order of magnitude beyond the hardest completed word at $e=2$, and is
recorded as inconclusive. Engine
verification with a corrected instance generator: all fifteen column words at $(2,5)$ are
exhausted, including the word this paper records below as surviving every filter.
\end{remark}

\begin{remark}[The gap is exactly one crossing, and there is no cheaper repair]\label{O-n1gapexact}

The step admits a sharp reduction. Let $Q$ be the tile across $T_k$'s $b$-edge
$[V_k,\,((k+1)c,0)]$, whose base end lies on $\partial ABC$ and whose end $V_k$ is interior. Measured
from the anchored base end, $Q$ lays $a$, $b$ or $c$ along that edge:
\begin{itemize}
\item $b$: an exact match, so $Q$ is the partner $P_k$ --- the case the induction wants;
\item $c$: overruns $V_k$ by $c-b=e^2$;
\item $a$: falls short of $V_k$ by $b-a$, leaving a run of that length still to be covered.
\end{itemize}
The third case does not close either, because $b-a$ is representable in $\langle a,b,c\rangle$ for no
member whatever. Indeed if $xa+yb+zc=b-a$ then $y=0$, since $y\ge1$ would give $yb\ge b>b-a$;
reducing $x\,ef+z\,f^2=f^2-e^2-ef$ modulo $f$ then gives $f\mid e^2$, hence $f=1$ by coprimality.
(\texttt{gap\_b\_sub\_a}; checked on all $29435$ coprime members with $e\le200$, $e<f<4e+3$, with no
representation in any of them.)

So both surviving branches force a tile edge past $V_k$, and the induction step of
Theorem~\ref{C-thm:n1} is \emph{equivalent} to the assertion that nothing overruns $V_k$. That is the
one-end-anchored pinning question in its purest form: the base end cannot be overrun, because the
continuation would leave the target, while $V_k$ is interior and, as computed in
Remark~\ref{C-rem:n1gap}, is not blocked by any tile already forced. Repairing Theorem~\ref{C-thm:n1} at
$e\ge2$ is therefore neither more nor less than settling that question, and the same question is what
stops the corner cascade of Remark~\ref{C-rem:route1uniform} and the filler dichotomy of
Section~\ref{C-sub:sidep}. Three routes, one wall.
\end{remark}

\begin{remark}[Counting cannot decide the crossing question]\label{O-censusblind}

Since an overrun at $Y$ is equivalent to the tiles with a \emph{vertex} at $Y$ forming a
$\pi$-figure, it is natural to attack it with the vertex census of
Lemma~\ref{C-lem:census}, which counts exactly those figures. That cannot work, for a structural
reason worth recording so it is not attempted again.

Write $n_1$ for the number of $\{\gamma,\alpha,\beta\}$ figures and $n_2$ for the number of
$\{3\alpha,2\beta\}$ figures, and $v_1,\dots,v_4$ for the four interior types. The census consists of
one equation per angle, and its identity is obtained by subtracting the $\beta$ equation from the
$\alpha$ equation:
\[
v_1 \;=\; 1+n_2+v_3+2v_4 ,
\]
in which $n_1$ \emph{cancels}. The third equation then pins $n_1=N-3v_1-2v_2-v_3$, but supplies no
bound of its own. So the census constrains $\{3\alpha,2\beta\}$ and says nothing bounding about
$\{\gamma,\alpha,\beta\}$.

That is precisely the wrong way round for us: at $V_k$ the forced tiles are $T_k$, presenting
$\gamma$, and $P_{k-1}$, presenting $\beta$, and $\{\gamma,\alpha,\beta\}$ is the only $\pi$-figure
containing both. The census is therefore blind to exactly the configuration in question, and no
counting argument built on it can decide the crossing.

The kernel-checked $99$-tiling illustrates both points. Its vertex figures are: on the boundary,
$22$ of type $\{\gamma,\alpha,\beta\}$, with the corners contributing $\{\beta\}$ twice and
$\{3\alpha\}$ once; in the interior, $50$ of type $\{\gamma,\alpha,\beta\}$, $2$ of type
$\{3\alpha,2\beta\}$, and $(v_1,v_2,v_3,v_4)=(4,7,1,0)$. The identity reads $4=1+2+1+0$, exactly. And
the $50+2=52$ interior $\pi$-figures agree exactly with the $52$ edges carrying an interior vertex
found in Remark~\ref{C-rem:noedgetoedge} by an unrelated computation --- a useful cross-check on both.
Note the ratio: $50$ of the $52$ crossings are of the type the census cannot see.
\end{remark}

\begin{remark}[What this buys, and what it does not]\label{O-ptwoscope}

Corollary~\ref{C-cor:ptwo} is independent of Theorem~\ref{C-thm:n1} and of the base walk: it
constrains the side at any member. On the close pairs
with $f>2e$ --- $203$ of them with $e\le40$, the smallest being $(3,7)$ with $N=138$, then $(4,9)$
with $N=227$ and $(5,11)$ with $N=338$ --- the side parameter is bounded by $2$. Checked on all
$126003$ close pairs with $f>2e$ and $e\le1000$: $\lfloor(f-1)/e\rfloor=2$ in every one.

On that family Corollary~\ref{C-cor:basedi2e} also reduces the base to two columns. It does not reduce
it to one: that step needs Theorem~\ref{C-thm:n1} at $e\ge2$, which Remark~\ref{C-rem:n1gap} leaves open.
So for these members the current state is two base columns and two side parameters, not one of each.

The contrast with $e=1$ is the point. There the range is $1\le p\le f-1$ and grows without bound
with $f$, which is why the corner cascade of Conjecture~\ref{C-conj:advance} has to be run as an
induction over block lengths. Here the range is $\{1,2\}$ for every member of an infinite family, so
the side question is a fixed finite case analysis rather than an induction. This is the first
sub-family in which that happens.

At $(3,7)$, concretely: $p=1$ gives a side of $7$ $a$-edges and $4$ $c$-edges, and $p=2$ gives $14$
$a$-edges and $1$ $c$-edge. Excluding those two remains open; the reduction bounds the problem, it
does not solve it.
\end{remark}

\begin{remark}[Why this does not exclude a member, and the trap it sets]\label{O-parityscope}

It is tempting to combine Lemma~\ref{C-lem:parity} with the boundary walks. On a member all of whose
walks have odd edge-counts --- $(3,7)$ is one, with $k=4p+7$ on a side and $9$ or $13$ on the base
--- the boundary contributes $\sum k_i-3$ straight figures, an even number, against the odd value
$N+1=139$ that Lemma~\ref{C-lem:parity} demands. That looks like an exclusion. It is not.

$S$ counts \emph{every} straight figure, and a straight figure need not lie on the boundary: at an
interior $T$-junction one tile has the vertex in the relative interior of a side, contributing a
straight angle and no corner, while the tiles on the other side have corners summing to $\pi$. Such a
vertex is a $\{\gamma,\alpha,\beta\}$ or $\{3\alpha,2\beta\}$ figure and is counted in $n_1+n_2$.
Writing $T$ for their number, $S=\sum k_i-3+T$, so the parity constrains only $T$: at $(3,7)$ every
tiling has an \emph{odd} number of interior $T$-junctions. Since non-edge-to-edge incidences are
exactly what this problem must allow, nothing is excluded.

This is the failure mode Lemma~\ref{C-lem:census} already warns of --- corner counting cannot see the
side structure --- and the boundary walks are side structure. It is recorded here because the
temptation is sharp and the arithmetic is otherwise correct.
\end{remark}

\begin{remark}[What this does and does not extend]\label{O-abnecessary}

Theorem~\ref{C-thm:secondc} uses $e^{2}+ef<f^{2}$ at exactly one step: that $Y$'s edge from $P$ towards
$J$, being of length at most $|PJ|=a$, is therefore \emph{equal} to $a$. That inference needs $a$ to be
the shortest side, and it fails without it: at $(2,3)$, $(3,4)$, $(4,5)$, $(5,6)$ the tiles are
$(6,5,9)$, $(12,7,16)$, $(20,9,25)$, $(30,11,36)$, each with $b<a$, so a $b$-edge also fits inside
$|PJ|$; since $b$ is a flank of $\alpha$, the conclusion ``$Y$ does not present $\alpha$'' no longer
follows.

Proposition~\ref{C-prop:nogolden} removes the hypothesis from Theorem~\ref{C-thm:secondc}, and with it
from Corollary~\ref{C-cor:pbound}: the bound $pe+2\le f$ holds at every coprime member. It does
\emph{not} remove it from Theorem~\ref{C-thm:nobothmirror}, whose third cross product is a positive
multiple of $b^{2}-a^{2}$ and changes sign at the threshold; there the hypothesis is genuinely part of
the statement.
\end{remark}

\begin{remark}[Scope, and the status of this theorem]\label{O-efminus1}

The family $e=f-1$ was outside the reach of \S\ref{C-sub:apex} until Proposition~\ref{C-prop:nogolden}
removed the golden-ratio hypothesis: $f/(f-1)<\varphi$ for every $f\ge3$, so every member of it has
$b<a$. Its tile counts are $N=2f^{2}+2f-1$, giving $11, 23, 39, 59, 83, 111, 143, 179, 219, 263,
311,\dots$, of which $23$ and $39$ were previously settled by exhaustive search, $11$ by
Theorem~\ref{C-thm:basebeta-e1}, and $59$ by the frontier results for $N\le80$. Among the remainder the
base-$\beta$ candidates, that is the primes congruent to $11$ modulo $12$, are
\[
83,\;179,\;263,\;311,\;419,\;479,\;683,\;839,\;1103,\;1511,\;2111,\;2243,\;2663,\;2963,\;3119
\]
for $f\le39$. Whether the family contains infinitely many such primes is a Bunyakovsky-type question
and is not claimed here.

The theorem discharges Hypothesis~\ref{C-hyp:walls} at these members; it does not by itself exclude the
corresponding integers, which requires the remaining branches.

We record the label honestly. The chain is PROVED, not VERIFIED. Lean carries the arithmetic
(\texttt{ApexRigidity.b\_gt\_f}, \texttt{.eb\_gt\_a}, \texttt{.e\_ge\_two\_of\_b\_lt\_a},
\texttt{.side\_p\_bound}), the area accounting (\texttt{.area\_above\_chord}, \texttt{.middle\_fraction})
and the vertex budget of Corollary~\ref{C-cor:noTP} in both branches
(\texttt{SecondEdge.noTP\_budget}, \texttt{.noTP\_remainders}, \texttt{.T2\_not\_alpha\_at\_P}). What is
not formalized is the planar geometry underneath: that $T_1$ presents $\gamma$ at $P$ and that $T_2$'s
chord trace ends there, which are the hypotheses those statements take. The non-representability of $a-b$
was checked symbolically, with $(f-2)(2f-1)-(a-b)=(f-1)^{2}$, and by exhaustive enumeration for
$3\le f\le39$. Two adversarial passes have been made, both within the session that produced the proof. The first
found a genuine defect --- Corollary~\ref{C-cor:pbound} was stated under a hypothesis that fails at
$e=f-1$, the very members it was invoked for --- now repaired by restating that corollary under
$\gcd(e,f)=1$. The second checked the side equation, integrality of $p$, the case $p\ge2$ including
$f=2$, that $J$ is interior to the side so that the tile $Y$ exists at all, the edge counts, absence of
circularity, and that Theorem~\ref{C-thm:apexconfig} and Corollary~\ref{C-cor:noTP} use nothing beyond
family-wide facts; it found no further defect.

The identification with Hypothesis~\ref{C-hyp:walls} is the one discussed in
Remark~\ref{C-rem:wallsfix} and is not settled; the theorem is stated in terms of $p$ for that reason. Theorem~\ref{C-thm:efminus1} is therefore unconditional as a statement about $p$, and
inherits that identification when read as a statement about the hypothesis.

A second point in its favour: since $N=m^{2}(3f^{2}-e^{2})$, a prime $N$ forces $m=1$, so at each of
the fifteen candidates the branch has no other scale to consider.

Neither pass is what Rule 6 asks for, namely a review in a fresh session with no memory of the
derivation, and that remains outstanding. Each of the fifteen candidate integers above has, we note, a
unique base-$\beta$ member, so discharging the hypothesis at $(f-1,f)$ discharges it for that
integer's only member of this branch.
\end{remark}

\begin{remark}[What this settles, and exactly where it stops]\label{O-apexscope}

Corollary~\ref{C-cor:pbound} is a second and independent proof of Theorem~\ref{C-thm:ptwodead} on the
tight subfamily, by a route that uses no chord bound and no height identity --- only the apex
identities and the boundary at $J$. It also completes the first bullet of
Remark~\ref{C-rem:dichotomy}: there the $a\,c$ ending was excluded only in the figure
$\{\gamma,\alpha,\beta\}$ and only with the $\alpha$ adjacent to $C$, whereas
Theorem~\ref{C-thm:secondc} excludes it in \emph{both} straight figures and for every $p$. Hence at
$(3,7)$ the side word of an equal side must end $c\,c$; the two cases $X=\beta$ and $X=\gamma$ of
Remark~\ref{C-rem:dichotomy} both persist, but now only inside that ending.

It does not reach $p=1$ at any new member, and the reason is arithmetic rather than incidental.
Killing $p=1$ by this bound needs $e\ge f-1$; with $e^{2}+ef<f^{2}$ that forces $f=e+1$ and then
$e^{2}<e+1$, so $e=1$ and $f=2$ --- the member already closed by
Theorem~\ref{C-thm:basebeta-e1} (\texttt{scope\_limitation}). The hypothesis $e^{2}+ef<f^{2}$ says
$f/e$ exceeds the golden ratio, and $e\ge f-1$ says it is close to $1$; the two cannot hold together
beyond $(1,2)$.

The obstruction to iterating is equally precise. Corollary~\ref{C-cor:noTP} needed \emph{two} tiles
presenting $\gamma$ at $P$, and the second, $T_2$, exists only because the apex figure is $3\alpha$.
At the analogous point one level down --- the far end of the next $c$-tile's $a$-edge, which does lie
exactly on the next chord, by the same identities --- only one $\gamma$ is forced. Supplying a second
is what an induction would require, and we do not have it.
\end{remark}

\begin{remark}[What this buys, and what it does not]\label{O-inflvalue}

This is a reformulation, not a new obstruction: the question ``is the inflation rigid'' and the
question ``is $p=0$'' are the same question, and at $(1,3)$ Theorem~\ref{C-thm:walls13} already answers
it. Reformulations earn their place only by new attack surface, and here there is one, namely size.
The inflation carries $f^2$ tiles against the target's $N=3f^2-e^2$:
\[
\begin{array}{rrr}
(e,f) & f^2 & N\\\hline
(1,3) & 9 & 26\\
(2,5) & 25 & 71\\
(3,7) & 49 & 138\\
(4,9) & 81 & 227
\end{array}
\]
uniformly about $35\%$. At $(3,7)$ that is a $49$-tile question in place of a $138$-tile one, with the
same answer for the west side. Proposition~\ref{C-prop:inflbdy} carries no unproved geometric hypothesis:
the $\gamma$-corner, the only one that can split, is not used.
\end{remark}

\begin{remark}[What this settles and what it does not]\label{O-inflrigidscope}

Remark~\ref{C-rem:sidenoa} draws exactly this inference --- ``the complete west block is the tile scaled
by $f$, so its far side \dots\ is subdivided into exactly $f$ $c$-edges'' --- and asserts the middle
step rather than proving it. Corollary~\ref{C-cor:sidenoa-proved} proves it at the three members where the
computation has been run, $(4,9)$ included.

It does not decide $p$ at $(3,7)$. The inflated tile is the west corner block only if that block is
complete, which is the west half of Hypothesis~\ref{C-hyp:walls} itself; the corollary is conditional on
exactly the hypothesis it serves. What has been removed is an unproved step \emph{inside} that
implication, not the implication's hypothesis.

At $(1,3)$ the theorem also reproves Theorem~\ref{C-thm:walls13} by an independent route, on $9$ tiles
rather than $26$, which is the regression check that the reformulation of \S\ref{C-sub:inflation}
recovers a known case before being trusted elsewhere.
\end{remark}

\begin{remark}[What this buys]\label{O-orientmono}

Two things, on different fronts.

On the inflation, The junction analysis of \S\ref{C-sub:inflation} (machine-checked as \texttt{junction\_three}) leaves three completions at each junction of the $k=2$ word
$c\,a^{f}c$, hence $3^{f-1}$ patterns a priori. Proposition~\ref{C-prop:orientmono} leaves $f+1$, one for
each value of $j$: the branching is linear in $f$, not exponential. That is consistent with the
measured search depth, which grows like $2.65f$ rather than with $f^2$.

On the thin family, nothing in the proof used the inflation --- only that an $a$-edge's two ends carry
$\beta$ and $\gamma$. So the proposition applies verbatim to an $a$-run on an \emph{equal side}, and in
particular to the runs enumerated in the second gap of Conjecture~\ref{C-conj:advance}. That enumeration
is over the double-$c$ walks $(fk,0,f-k)$ with $1\le k\le f-3$; the number of such words grows quickly
with $f$ \textup{(}$1,7,42,263,1823,\dots$ for $f=4,\dots,8$\textup{)}, so it is not a finite check, but
the orientation branching within each word is now linear rather than exponential. This is the first
tool bearing on that gap which is uniform in $f$; the cascade argument as written does not use it.
\end{remark}

\begin{remark}[Scope, and why this is not a proof of Theorem~\ref{C-thm:inflrigid}]\label{O-inflparity}

Two limitations, both essential.

First, $\gcd(e,f)=1$ forbids $e$ and $f$ both even, so ``$e\not\equiv f$'' is exactly ``one of them is
even''. The proposition is therefore \emph{silent} on the both-odd members --- which include $(1,3)$ and
$(3,7)$, i.e.\ every member the crux actually needs. It fires at $(2,5)$ and $(4,9)$, where
$3f^2-(4f-e)$ equals $57$ and $211$, and those are precisely the two members the search found cheap.
That correlation is the content of the proposition: it explains the observed cost profile, it does not
extend the theorem. The $(4,9)$ search has since returned \texttt{EXHAUSTED\_NO\_TILING} in
$5\,161\,190$ nodes, which agrees with the parity verdict there and is strictly stronger, the search
assuming no edge-to-edge condition.

Second, edge-to-edge is an extra hypothesis, and it is not a mild one here. The main paper does not
assume it, and the T-junction whose exclusion is the whole difficulty of \S\ref{C-sub:apex} is exactly an
edge-to-edge failure; assuming it would be assuming a neighbourhood of what is to be proved. This is the
same wall that Remark~\ref{M-rem:parityfail} hits: the $b$-edge pairing argument there fails precisely
because a longer edge may straddle a $b$-edge instead of abutting it, and both genuine tilings exhibit
such straddles. Proposition~\ref{C-prop:inflparity} is what survives once straddles are hypothesised away,
which is why it proves less than it appears to.

So the uniform proof of Theorem~\ref{C-thm:inflrigid} remains open. What Proposition~\ref{C-prop:inflparity}
shows is that any such proof must use more than counting on the both-odd members, since the counting
invariant is identically satisfied there. \textup{(}Machine-checked, \texttt{parity\_kills\_25},
\texttt{parity\_kills\_49}, \texttt{parity\_silent\_witnesses}.\textup{)}

One caveat should be entered in the hypothesis's favour. The edge multiset of the two standard
inflations was computed from the certified tilings in \texttt{code/engine/inflation/}: at $(1,3)$ every
edge has multiplicity $1$ or $2$ with $9$ of the former and $9$ of the latter, and at $(3,7)$ with $21$
and $63$; no edge has multiplicity above $2$, and $2\cdot 9+9=27=3\cdot 3^2$, $2\cdot 63+21=147=3\cdot
7^2$. So the inflation that does exist \emph{is} edge-to-edge, and the boundary count is $3f$ on the
nose. That is a positive control sensitive to exactly the quantity the proof above manipulates --- a
mis-stated $B$ would break it --- and it is recorded because the vertex-census parity of
Remark~\ref{C-rem:parityscope} was retracted after passing a control that could not see its defect.
What the control cannot do is speak for the $p=1$ boundary, since no such tiling exists to measure.
\end{remark}

\begin{remark}[The last junction, and what is still missing]\label{O-lastjunction}

At $J_{2f}$ the two boundary neighbours are the $a$-tile $A_{2f}$, presenting $\gamma$ with a
$b$-ray inward, and the apex $c$-tile, presenting $\beta$ with an \emph{$a$}-ray inward; the filler's
rays are $b$ and $c$. Since neither $b$ nor $c$ equals $a$, \emph{$J_{2f}$ cannot be edge-to-edge},
and each placement leaves a residual run: $c-a$, or $b-a$ together with $c-b=e^2$. All three are gaps
of $\langle a,b,c\rangle$ --- $b-a$ by \texttt{gap\_b\_sub\_a}, $e^2$ by \texttt{gap\_e\_squared},
and $c-a$ by \texttt{gap\_c\_sub\_a\_tight}, which holds precisely for $e\ge2$ since $c-a=2a$ at
$(1,3)$. Moreover $e^2<a$, so no tile edge fits inside that run at all. Hence the covering
\emph{must} overrun a filler vertex.

Each placement therefore forces the covering to overrun a filler vertex. That alone is not an
exclusion, and the exclusion comes instead from a chord bound, which we now state.
\end{remark}

\begin{remark}[The audit of Theorem~\ref{C-thm:ptwodead} against the apex rigidity]\label{O-ptwoaudit}

Theorem~\ref{C-thm:ptwodead} predates \S\ref{C-sub:apex}, and the two overlap. Reconciling them shows the
earlier enumeration was doing more work than necessary, and we record this rather than leave two
arguments whose relation is unclear.

Remark~\ref{C-rem:lastjunction} has the $\alpha$-filler standing between the apex $c$-tile and the
$a$-tile, so it lays one of its flanks --- a $b$ or a $c$ --- along the apex tile's $a$-ray. Both
exceed $a=|JP|$, so $P$ would lie \emph{interior} to that edge, contributing $\pi$ at $P$. But
Theorem~\ref{C-thm:apexconfig} puts $T_1$ at $\gamma$ there, and $T_2$'s trace on the chord ends at $P$,
so $T_2$ contributes at least $\gamma$ as well. The total is at least $2\gamma+\pi=2\pi+\alpha>2\pi$.

So that configuration does not occur at all, and the placements enumerated in the proof of
Theorem~\ref{C-thm:ptwodead} are placements of a tile that cannot be there. The conclusion is unharmed
--- indeed Corollary~\ref{C-cor:pbound} re-derives it, on a wider range of members and without any chord
estimate --- and Theorem~\ref{C-thm:lastjunction}'s conclusion, that the adjacent tile lays an $a$-edge,
is likewise immediate from the same count. What the audit removes is the impression that the tight
subfamily needs the height identity: it does not.
\end{remark}

\begin{remark}[Scope, and where $e\ge2$ is easier]\label{O-ptwoscopeb}

Theorem~\ref{C-thm:ptwodead} closes $p=2$ on the whole subfamily $f=2e+1$, hence at $(1,3)$, $(2,5)$,
$(3,7)$, $(4,9)$, $(5,11)$ and onward; verified independently by exact enumeration of the six
placements for every member with $e\le400$. Combined with Corollary~\ref{C-cor:ptwo} it leaves $p=1$ as
the only obstruction to Hypothesis~\ref{C-hyp:walls} on that subfamily.

It does not extend beyond $f=2e+1$: the height identity used above holds exactly there.

Two by-products. First, $c-a$ is a gap of $\langle a,b,c\rangle$ precisely for $e\ge2$, failing at
$(1,3)$ where $c-a=2a$; this is the first point in the problem where $e\ge2$ is easier than $e=1$.
Second, at $(3,7)$ the base walk $(0,e,2e)$ carries no $a$-edge, so the corner tile lays $c$ on the
base and $a$ on the side, forcing $p\ge1$ on both equal sides; with $p=2$ excluded,
Hypothesis~\ref{C-hyp:walls} at $(3,7)$ therefore requires the base walk $(f,e,e)$.\end{remark}

\begin{remark}[What the two no-go results mean]\label{O-nogoscope}

Proposition~\ref{prop:nogoauto} shows that the local combinatorics of the side is almost unconstrained:
the automaton forbids one adjacency and leaves $10\,560$ configurations. Proposition~\ref{prop:nogocensus}
shows that adjoining the census adds nothing, because the two counting equations are dependent. This
is the third independent confirmation of the warning in Lemma~\ref{C-lem:census} that corner counting
cannot see $p$, and it is also the exact reason the retraction described in
Remark~\ref{C-rem:parityscope} was possible: a counting argument at $(3,7)$ has one relation to spend
and it is spent.

The conclusion is directional, not merely negative. Any mechanism that closes $p=1$ must use metric
information --- positions, not multiplicities. Proposition~\ref{prop:straddle} is the first such fact.
\end{remark}

\begin{remark}[The wedge-filler hypothesis, and where it fails]\label{O-alphamin}

Several steps of the corner chain fill a wedge of angle exactly $\alpha$ with a single
$\alpha$-tile, justified in Lemma~\ref{C-lem:wallclimb} by ``$\alpha$ is the minimal angle''. That is
not the load-bearing fact. What the step needs is that no two smaller angles fit, that is
$2\beta>\alpha$, which with $3\alpha+2\beta=\pi$ reads
\[
\alpha<\tfrac\pi4 .
\]
The two conditions are different, and both directions of the difference occur. At $(e,f)=(3,4)$ one
has $b<a$, so $\beta$ is the minimal angle, yet $\alpha<\pi/4$ still holds and the step survives.

The criterion is exact and integral. Since $b^2+c^2-a^2=(f^2-e^2)(2f^2-e^2)$ and
$2bc=2f^2(f^2-e^2)$,
\[
\cos\alpha=\frac{2f^2-e^2}{2f^2}=1-\frac{e^2}{2f^2},
\]
which at $e=1$ is the familiar $(2f^2-1)/(2f^2)$. So $\alpha<\pi/4$ iff $2f^2-e^2>\sqrt2\,f^2$, and
squaring (both sides positive, as $f>e$) gives
\[
e^4-4e^2f^2+2f^4>0 .
\]
The separation condition $f^2>2ef+e^2$ implies it (\texttt{separated\_wedge\_ok}): from
$f^2-e^2>2ef$ one gets $f^4+e^4>6e^2f^2$ on squaring, and the claim follows with room to spare. So
the step is safe at every separated member, and every failure is a close pair. Among $e\le11$,
$f\le19$ the failures are
\[
(4,5),\ (5,6),\ (6,7),\ (7,8),\ (7,9),\ (8,9),\ (9,10),\ (9,11),\ (10,11),\ (10,13),\ (11,12),\ (11,13),
\]
including $(5,6)$, which is $N=83$. No claim in either paper is affected: Lemma~\ref{C-lem:wallclimb}
is invoked only at $e=1$, where $a=f<f^2-1=b$ and $\alpha<\pi/4$ both hold, and
Proposition~\ref{prop:unify} lists four ingredients not including it. What the criterion supplies is
a precise location for the obstruction: close pairs differ from separated members geometrically, not
merely in having weaker arithmetic filters, and an argument aimed at them must either avoid the
wedge-filler step or treat $e^4-4e^2f^2+2f^4\le0$ separately.

Two arithmetic strengthenings do \emph{not} help there, which is why the bound must be interior.
Reducing the base equation modulo $e$ gives $y+z\equiv0\pmod e$ by coprimality; and the
$\gamma$-injection is exactly a bipartite matching of the $a$- and $b$-edges to the junctions, whose
sharp criterion is Hall's condition rather than the counting bound. Neither removes any of the $54$
columns that survive the three unconditional filters over $e\le9$, $f\le15$: the first is implied by
the base equation, and the second collapses to the $\gamma$-trap $z\ge1$, since with $z\ge1$ a
matching always exists.
\end{remark}

\begin{remark}[What reach buys, and what it does not]\label{O-reachlaw}

The four kills at reach $R$ empty the window exactly for $f\le R+2$: reach $3$ gave $f\le5$,
reach $4$ gives $f\le6$, and each further row would give one further member. Theorem
\ref{C-thm:strippbound} raises the reach by one and no more, because the clearance bound it uses is
the fixed constant $T/b>2$. A proof for all $f$ therefore requires killing straddlers at
\emph{every} row, not at the first three. The natural route is that a ladder descends into strips
the chain has already forced, where every line of its direction is a shared $b$-edge chain broken
at the strip junctions on both sides at once --- unstaggered --- while a ladder is staggered by
$c-b=1$; a ladder in forced territory would have to run along such a line or cross tile interiors,
and both are impossible. What is missing is the bookkeeping: that the region a ladder descends
through, roughly one half-column east per strip, lies inside the forced staircase at the moment
the fork is being decided.
\end{remark}

\begin{remark}[What this does to the induction]\label{O-e2induction}

The collar of Theorem~\ref{C-thm:realize12} runs on unit columns, so it is unavailable at $e\ge2$;
only the wide columns of Proposition~\ref{C-prop:widecol} remain, and they advance $m$ in steps of
$f$. Consequently the clean law ``tileable iff $m\ne1$'' proved for $(1,2)$ need not persist:
the spectrum may genuinely depend on $e$. For $(2,3)$, where $f=3$, the subdivision closure and
the wide-column law together would give every $m\ge4$ and every even $m$ from a single seed
$\Delta_2$, leaving exactly $m=3$ unreached from below; so the decisive experiment for the whole
$e\ge2$ story is the realizability of $\Delta_2$ at $(2,3)$, i.e.\ $N=92$. That search is
running; we record the reduction rather than guess its outcome.
\end{remark}

\begin{remark}[The chord programme is exhausted]\label{O-chordexhausted}

It is natural to hope that the family of chords does collectively what no single chord does: there
are $37$ admissible heights $\rho=n/49$ with $n=3i+7j$, each carrying its own non-integral area
$138\rho^{2}$, and a tile lying below the chord at $J$ lies above the chord at some lower junction. We
record that this hope is also unfounded, so that the reader need not pursue it.

For a tile $T$ and a chord $C$ write $\varphi_T(C)\in[0,1]$ for the fraction of $T$'s area lying above
$C$; thus $T$ straddles $C$ exactly when $\varphi_T(C)\in(0,1)$. The area statement at height $\rho$ is
\[
\sum_{T}\varphi_T(C_\rho)\;=\;138\rho^{2},
\]
and this is nothing but additivity of area over the tiling: the left-hand side \emph{is} the area
above $C_\rho$, decomposed tile by tile. It therefore holds identically, for every tiling of the
target, and the $37$ admissible heights supply $37$ instantiations of one identity rather than $37$
constraints.

The only information extractable from the family is the integrality bit --- where $138\rho^{2}\notin
\mathbb Z$, some $\varphi_T(C_\rho)$ lies strictly between $0$ and $1$ --- and that is precisely
Proposition~\ref{prop:straddle}. The nesting adds nothing further: the set of heights straddled by a
given tile is an interval, and $138$ tiles cover $37$ heights with room to spare.

This is the same degeneracy as Proposition~\ref{prop:nogocensus}, in a metric rather than a
combinatorial register: a single global identity, instantiated many times, cannot separate tilings.
Taken together with Remark~\ref{rem:chordscope} it closes the chord programme. What the chord method
yields at $(3,7)$ is exactly Propositions~\ref{prop:straddle} and \ref{prop:chorddecomp}, and closing
$p=1$ requires a mechanism that is neither counting nor chord-area. We do not have one, and we prefer
to say so.
\end{remark}

\begin{remark}[Where the chain stops being free]\label{O-chainobstruction}

Lemma~\ref{lem:cornerstep} iterates only as far as the same two ingredients survive. The fill of
$\alpha+\beta$ is unique at every member, so the two tiles are always an $\alpha$ and a $\beta$;
what decides whether the next base edge is confined to $\{a,c\}$ is their angular order, and that
is settled by whether the $\alpha$-tile is the partner across the previous tile's whole $b$-edge.
At the corner this is unconditional, because the $b$-chord runs from a base vertex to a side
vertex. At the $k$-th step the chord's upper end is interior, and pinning it is precisely the
crossing question of Remark~\ref{C-rem:straddle}. Both remaining cases, the residual configurations
at $e=1$ and the exclusion of $1\le p\le(f-1)/e$ at $e\ge2$, reduce to that single question.

That question is not answered by the hypothesis $m=1$ alone. Instrumenting the certified
exhaustions to record every partial configuration they build shows the overshoot occurring
freely: at $(1,2)$ it appears in $31$ of the $52$ recorded frontiers and at $(1,3)$ in $1128$,
the first already at depth $3$ with three tiles placed, a $c$-edge from $(2,0)$ to
$\bigl(\tfrac{11}2,\tfrac{\sqrt{15}}2\bigr)$ covering a $b$-edge from $(2,0)$ and overshooting
it by $c-b=1$. Every such configuration is eventually refuted, but not locally: the partner
pinning cannot be established by excluding the overshoot near the chord.
\end{remark}

\begin{remark}[The thick program]\label{O-thickprogram}

The three verified facts above change the character of the $e\ge2$ case. The chain machinery
ports with strictly stronger tools: the wall climb and the first-run chirality are
member-independent (their proofs use only the angle identities and the boundary overshoot
$c-b=e^2>0$), and Lemma~\ref{C-lem:monochotomy} removes the $c$-chord fork entirely, so the thick
chain branches strictly less than the thin one. What remains is the finite bookkeeping: the base
letters run through a trichotomy rather than a free word, the side violations are quantized to
$p\in f\mathbb{N}$, and the reach analysis of the thin case applies verbatim where the blocks are
short. The $(2,3)$, $(2,5)$ and $(3,4)$ members are certified by full exhaustions ($19{,}677$,
$553{,}417$ and $825{,}724$ nodes; the second is the prime $N=71\equiv11\pmod{12}$, the third a
close pair). The close pair refutes deeper ($8.3f$ against $5.7f$ and $5.2f$), consistent with
its additional base decompositions.
\end{remark}

\begin{remark}[Why this is the right direction]\label{O-straddlescope}

At the last junction, $n=7$ and the area above the chord is $138/49=2.816\ldots$ tiles, while the
chord has length $414/7=59.14\ldots$ and a tile meets a line in a segment of length at most $c=49$;
so at least two tiles meet that chord, and at least one straddles it. This is exactly the region
below the chord that the enumeration of Theorem~\ref{C-thm:ptwodead} does not reach --- see the
discussion of the $c\,c$ ending in Remark~\ref{C-rem:dichotomy} --- and it is the reason the surviving
configurations survive: not that the chord bound fails there, but that it was never applied there.
Whether the straddling constraint can be sharpened into a contradiction is open.
\end{remark}

\begin{remark}[What the decomposition gives, and the rung that remains]\label{O-chordscope}

Proposition~\ref{prop:chorddecomp} is the first statement that constrains the configuration
\emph{below} the chord at $J$, the region Remark~\ref{C-rem:dichotomy} identifies as untouched. It
reduces the possibilities there to a finite list: a flush total in $\{0,21,40,42,49\}$ together with a
straddling remainder of $59.142\ldots$, $38.142\ldots$, $19.142\ldots$, $17.142\ldots$ or
$10.142\ldots$, distributed over at least one and at most a bounded number of straddling tiles.

What it does not give is rigidity: the cross-section of a straddling tile varies continuously with its
position, so no individual value is forced. The natural hope is that the $\gamma$-grading of
\S\ref{C-sub:gammagrading} restores discreteness --- a straddler's edges have directions that are
integer multiples of $\gamma$ modulo $\pi$, so one might expect its cross-section to be drawn from a
countable set fixed by its label. \emph{That hope is unfounded}, and we record why, since it is the
first thing a reader will try.

The grading constrains directions, not positions. Tile vertices lie in the $\mathbb Z$-module
generated by the admissible edge vectors, and for a horizontal chord what matters is the projection
of that module to the vertical axis. Since $\gamma=(\pi+\alpha)/2$ with $\alpha/\pi$ irrational, that
projection is generated by the values $\ell\sin(k\gamma)$ for $\ell\in\{21,40,49\}$ and $k$ in the label
window, and these are rationally independent: the subgroup they generate is dense in $\mathbb R$.
So no cross-section is forced, and the chord decomposition cannot by itself close $p=1$.

What the decomposition does establish is that the obstruction is located strictly \emph{below} the
chord. Above it there is only $138/49=2.816\ldots$ of a tile's area, so that region is enumerable up
to boundedly many tiles --- which is precisely the enumeration Theorem~\ref{C-thm:ptwodead} performs.
The surviving configurations of Remark~\ref{C-rem:dichotomy} place the decisive tile below the chord,
where no area bound confines it. Closing $p=1$ therefore requires a constraint that reaches downward,
and we do not have one.
\end{remark}

\begin{remark}[The residue at $f\ge6$]\label{O-f6residue}

For $f\ge6$ the surviving words form an explicit thin set, symmetric under reversal; at $f=6$ it
is a single class: $b$ at position $5$, $c$ at position $7$ ($a\,a\,a\,a\,b\,a\,c\,a$), read from
a corner whose block is at least $3$. The blocking structure is always the same: the chain's
column-$3$ line needs a row-$3$ brick, the row-$3$ fork's $A$-branch lays its south cover on the
row-$2$ floor, and the floor is edged only as far as the staircase has reached --- which the $b$
itself interrupts. Closing the row-$3$ fork is the entire remaining content of
Conjecture~\ref{C-conj:advance} at $e=1$.
\end{remark}

\section{Routes closed, with their failure points}\label{sec:dead}

The following approaches to this branch are closed, and are recorded so that they are not retried.

\begin{enumerate}
\item \textbf{The whole class of linear direction invariants.} Weighting a directed edge by a
function of its direction alone, antisymmetric under $\theta\mapsto\theta+\pi$, yields a
one-parameter family of identities; every member is a specialisation of the master identity of
\cite[Rem.~4.9]{Main}. The associated ideal-membership condition is satisfied identically for every
root of unity and every $(e,f)$, and the optimal $\ell^1$ bound the family can produce is
$m(3f-2e)(f+e)/f$, which is smaller than $N$ by a factor at least $11/6$ and growing like $mf$. So
no choice of weight, kernel, moment ladder or duality can exclude a candidate. Both statements are
machine-checked in \texttt{lean/LambdaFactor.lean}.
\item \textbf{Closing $e=1$ by proving $R\ge2$ on every side.} False at the base: the walk with
$R_{\mathrm{base}}=2$ is exactly the one the corner parallelogram kills, and the survivor has
$R_{\mathrm{base}}=1$. On the base, ``$R\ge2$'' restates the remaining problem rather than reducing
it. It does hold on the two equal sides, and is proved there.
\item \textbf{Boundary tile counts.} The number of tiles meeting $\partial ABC$ grows like the
number of boundary edges, and on the genuine certificates it is $28$ of $44$ and $44$ of $99$ --- a
falling fraction. Comparing it against $N$ therefore never yields a contradiction.
\item \textbf{The far side of the mismatch ray}, under the hypothesis that the ray is a complete
wall: that side is the tile scaled by $f$, and is tileable. No obstruction can come from there.
\item \textbf{Local combinatorics of an equal side, with or without the census}, for $p=1$ at
$(3,7)$. The junction automaton forbids a single adjacency and leaves $10\,560$ configurations
(Proposition~\ref{prop:nogoauto}); adjoining the corner census adds nothing, its two counting
equations being dependent (Proposition~\ref{prop:nogocensus}). Together with
Lemma~\ref{C-lem:census} this is the third confirmation that corner counting cannot see $p$.
\item \textbf{The chord/area programme}, beyond Propositions~\ref{prop:straddle}
and~\ref{prop:chorddecomp}. Two continuations both fail: the $\gamma$-grading does not discretize
straddle positions, because it constrains directions and the vertical projection of the vertex module
is dense (Remark~\ref{rem:chordscope}); and the family of $37$ admissible chord heights supplies $37$
instantiations of one identity --- additivity of area --- rather than $37$ constraints
(Remark~\ref{rem:chordexhausted}). Closing $p=1$ needs a mechanism that is neither counting nor
chord-area.
\item \textbf{Boundary-word invariants (Conway--Lagarias tiling groups)}~\cite{ConwayLagarias}. These
are the natural non-abelian strengthening of the direction invariants of item~1, and are known to be
strictly stronger than colouring and weighting arguments --- which is exactly the class item~1 kills
here. They nevertheless do not apply, for two independent reasons. First, the method is formulated
for regions assembled from cells of a regular lattice, and the vertices of a tiling of our target lie
in no lattice: the edge vectors have lengths in $\{21,40,49\}$ and directions $k\gamma$ with
$\gamma=(\pi+\alpha)/2$ and $\alpha/\pi$ irrational, so the vertical projection of the
$\mathbb Z$-module they generate is dense in $\mathbb R$. Second, tilings here need not be
edge-to-edge, so the boundary word of a region is not determined by its combinatorics.
\end{enumerate}



\section{Two mechanisms that provably cannot close \texorpdfstring{$p=1$}{p=1}}\label{sec:nogo}

Remark~\ref{C-rem:dichotomy} ends by saying that closing $p=1$ ``requires a further mechanism''. That
is a hedge, and it can be replaced by theorems. This subsection records two natural mechanisms and
proves that neither can work, and one arithmetic fact that survives and points the other way. The
member is $(e,f)=(3,7)$ throughout: tile $(a,b,c)=(21,40,49)$, target $(343,343,414)$, $N=138$.

Fix an equal side with $p=1$. Its word is a permutation of $a^{7}c^{4}$ ending in $c$
($7\cdot21+4\cdot49=343$), so it has $11$ edges and $10$ interior junctions. Each edge is laid by a
tile, and the two angles of that tile at the ends of the edge are the two angles \emph{not} opposite
it: an $a$-edge is flanked by $\beta$ and $\gamma$, a $c$-edge by $\alpha$ and $\beta$. Call the
choice of which flank sits at the lower end the \emph{orientation} of the edge.

\begin{lemma}[Endpoint angles]\label{lem:endpoints}\lab{PROVED}
The angle at the top of the last edge is $\alpha$, and the angle at the bottom of the first edge is
$\beta$.
\end{lemma}

\begin{proof}
By Lemma~\ref{C-lem:census} the apex fills uniquely as $\{3\alpha\}$ and each base corner uniquely as
$\{\beta\}$. The last edge of the side is a $c$-edge, whose flanking angles are $\alpha$ and $\beta$;
the tile at the apex presents $\alpha$, so the top angle is $\alpha$. The single tile at the base
corner presents $\beta$, so the bottom angle of the first edge is $\beta$.
\end{proof}

\begin{proposition}[The junction automaton is consistent, and hugely so]\label{prop:nogoauto}\lab{PROVED}
Regard the side as a transfer system: a state is the angle at the top of the current edge, and a
transition is admissible when the pair (top angle of one edge, bottom angle of the next) extends to a
straight figure at their junction. Then the only inadmissible pair is $(\gamma,\gamma)$, which can
occur only at a junction between two $a$-edges. Consequently, over all $120$ words $a^{7}c^{4}$
ending in $c$ and all $2^{11}$ orientations, exactly $10\,560$ configurations satisfy every junction
constraint together with Lemma~\ref{lem:endpoints}.
\end{proposition}

\begin{proof}
A junction of an equal side carries a straight figure, so by Lemma~\ref{C-lem:census} it is
$\{\gamma,\alpha,\beta\}$ or $\{3\alpha,2\beta\}$. The figure $\{\gamma,\alpha,\beta\}$ has one angle
of each kind, so it admits a pair of side angles precisely when the two are distinct. The figure
$\{3\alpha,2\beta\}$ contains no $\gamma$, so it admits precisely the pairs drawn from
$\{\alpha,\beta\}$, all four of which leave an admissible remainder ($3\alpha$, $2\alpha+\beta$,
$2\alpha+\beta$, $\alpha+2\beta$ respectively). A pair is therefore inadmissible exactly when it is
equal and involves $\gamma$; and $\gamma$ occurs only as a flank of an $a$-edge. The count is a
finite check over $120\cdot2^{11}$ cases.
\end{proof}

\begin{proposition}[The census contributes one relation, and it is already spent]\label{prop:nogocensus}\lab{PROVED}
Let $x$ be the number of junctions of the side carrying $\{\gamma,\alpha,\beta\}$. Then $x\ge7$, and
the $\alpha$- and $\beta$-counts along the side are linearly dependent: they sum to the identity
$13=13$. Both $x=7$ and $x=10$ occur in solutions, so the census does not determine $x$.
\end{proposition}

\begin{proof}
The $7$ $a$-edges contribute one $\beta$ and one $\gamma$ each, the $4$ $c$-edges one $\alpha$ and
one $\beta$ each: $7\gamma$, $11\beta$, $4\alpha$ in all. By Lemma~\ref{lem:endpoints} the apex
absorbs one $\alpha$ and the base corner one $\beta$, leaving $7\gamma$, $10\beta$, $3\alpha$
distributed over the $20$ side-tile slots at the $10$ junctions. A $\{3\alpha,2\beta\}$ junction
contains no $\gamma$ and a $\{\gamma,\alpha,\beta\}$ junction exactly one, so the $7$ $\gamma$'s
occupy $7$ distinct junctions of the latter type, whence $x\ge7$.

Write $u$ for the number of junctions whose $\gamma$ lies on a side tile and whose other side slot is
$\beta$, and $(a_i,b_i)$ with $a_i+b_i=2$ for the side slots at the $10-x$ junctions of type
$\{3\alpha,2\beta\}$. Counting $\alpha$ and $\beta$ separately,
\[
(7-u)+(x-7)+\textstyle\sum_i a_i \;=\;3,
\qquad
u+(x-7)+\textstyle\sum_i b_i \;=\;10 .
\]
Adding and using $\sum_i(a_i+b_i)=2(10-x)$ gives $7+2x-14+20-2x=13$, an identity. The two equations
are therefore dependent. Taking $x=10$, $u=7$ satisfies both, as does $x=7$, $u=4+\sum_i a_i$ for any
$\sum_i a_i\le3$.
\end{proof}

\begin{remark}[What the two no-go results mean]\label{rem:nogoscope}
Recorded in \cite{Obstructions}. Full discussion in \cite{Obstructions}.
\end{remark}

\begin{proposition}[Every junction chord is straddled]\label{prop:straddle}\lab{PROVED}
Let $J$ be a junction of an equal side at distance $d$ from the apex, and let $C$ be the chord of the
target through $J$ parallel to the base. Then some tile of the tiling meets both open sides of $C$.
\end{proposition}

\begin{proof}
Since $d$ is a sum of edge lengths $21$ and $49$, we have $d=21i+49j$ and the height fraction is
$\rho=d/343=(3i+7j)/49$. Write $n=3i+7j$, so $1\le n\le48$ for a junction strictly between the apex
and the base. The chord cuts off a triangle similar to the target with ratio $\rho$, whose area in
units of the tile area is
\[
138\rho^{2}\;=\;\frac{138\,n^{2}}{2401}.
\]
Now $2401=7^{4}$ and $138=2\cdot3\cdot23$, so $\gcd(138,2401)=1$ and the quantity is an integer if
and only if $7^{4}\mid n^{2}$, i.e.\ $49\mid n$ --- impossible for $1\le n\le48$. Hence the area above
$C$ is not an integral number of tiles. If no tile met both open sides of $C$, every tile would lie
weakly on one side and the area above $C$ would be a sum of whole tile areas, a contradiction.
\end{proof}

\begin{proposition}[The general form]\label{prop:straddlegen}\lab{PROVED}
Let $(e,f)$ be any member with $\gcd(e,f)=1$. Then the area above the chord through any junction
strictly between the apex and the base is never an integral number of tile areas, so some tile meets
both open sides of that chord.
\end{proposition}

\begin{proof}
A junction at distance $d=ia+jc=f(ie+jf)$ from the apex has height fraction $\rho=m/f^{2}$ with
$m=ie+jf$ and $1\le m<f^{2}$, so the area above its chord is $Nm^{2}/f^{4}$ tile areas. Since
$\gcd(e,f)=1$ we have $\gcd(N,f)=\gcd(3f^{2}-e^{2},f)=\gcd(e^{2},f)=1$, so $f^{4}\mid Nm^{2}$ forces
$f^{4}\mid m^{2}$, hence $f^{2}\mid m$ prime by prime, which contradicts $m<f^{2}$. The rest is as in
Proposition~\ref{prop:straddle} (\texttt{ChordDecomp.area\_never\_integral}; the coprimality is
automatic on the tight subfamily, \texttt{ChordDecomp.coprime\_tight}). Checked by exhaustive
enumeration over all members with $e\le14$, $f\le31$ and all admissible junctions: $80\,299$ cases,
no exception.
\end{proof}

\begin{remark}[Why this is the right direction]\label{rem:straddlescope}
Recorded in \cite{Obstructions}. Full discussion in \cite{Obstructions}.
\end{remark}

\begin{proposition}[Chord decomposition at the last junction]\label{prop:chorddecomp}\lab{PROVED}
Let $C$ be the chord of the target through the last junction $J$ of an equal side at $(3,7)$, so
$|C|=414/7$. Let $\mathcal U$ be the set of tiles whose interiors meet the open half-plane above $C$.
Each $T\in\mathcal U$ either \emph{straddles} $C$ (its interior meets both open sides) or is
\emph{flush from above} (it lies weakly above the line of $C$). Then:
\begin{enumerate}
\item[(a)] a flush tile meets the line of $C$ in a full tile edge, of length $21$, $40$ or $49$;
\item[(b)] the traces on $C$ of the tiles of $\mathcal U$ have pairwise disjoint interiors and cover
$C$, so their lengths sum to $414/7$;
\item[(c)] at most two tiles of $\mathcal U$ are flush, and the multiset of flush lengths is one of
$\varnothing,\{21\},\{40\},\{49\},\{21,21\}$;
\item[(d)] at least one tile straddles, and if none is flush then at least two straddle.
\end{enumerate}
\end{proposition}

\begin{proof}
(a) For a tile lying weakly above the line of $C$, that line is a supporting line, so the intersection
is a face of the tile --- a vertex or an edge (\texttt{contact\_is\_edge}). A trace of positive length
excludes the vertex case, and the edge lengths are $a=21$, $b=40$, $c=49$.

(b) The closed region above $C$ is covered by the tiles, and a point of $C$ interior to the target is
a limit of points of that region, hence lies in a tile of $\mathcal U$, which is closed; so the traces
cover $C$. If two tiles of $\mathcal U$ had traces overlapping in a segment, both would occupy area
immediately above that segment, contradicting disjointness of interiors. Additivity of length gives
the sum.

(c) The flush lengths are integers drawn from $\{21,40,49\}$ and, by (b), sum to at most
$|C|=414/7=59.142\ldots$. The only multisets meeting that bound are the five listed: already
$21+40=61>|C|$, and $21+21+21=63>|C|$.

(d) By the computation in Proposition~\ref{prop:straddle} the area above $C$ is $138/49$ tiles, not an
integer, so a straddler exists. If no tile of $\mathcal U$ is flush, the straddling traces alone sum to
$|C|=59.142\ldots$; a segment contained in a tile has length at most the tile's diameter, which for a
triangle is its longest side $c=49$. Hence at least two straddlers.
\end{proof}

\begin{remark}[What the decomposition gives, and the rung that remains]\label{rem:chordscope}
Recorded in \cite{Obstructions}. Full discussion in \cite{Obstructions}.
\end{remark}

\begin{remark}[The chord programme is exhausted]\label{rem:chordexhausted}
Recorded in \cite{Obstructions}. Full discussion in \cite{Obstructions}.
\end{remark}








The clause ``that admits a solution'' in Proposition~\ref{C-prop:closepaircolumns} is not decided by
(iii): the bound there ignores the residue class of $x$, and over-permits. At $(e,f)=(8,17)$ it
allows $k=1$, while $8x+17z=47$ with $x\equiv8\pmod{17}$ has no solution with $x,z\ge1$, so no
column exists; the same happens at $551$ of the $1558$ close pairs with $e\le60$. Restoring the
residue class makes the test exact.


\begin{proof}
(i) By Proposition~\ref{C-prop:closepaircolumns}(i),(ii) a column satisfies $xe+zf=2ef-kb$ and
$x\equiv ke\pmod f$; conversely, for any $x$ in that class $f$ divides $2ef-kb-xe$, because
$xe\equiv ke^2$ and $2ef-kb\equiv ke^2 \pmod f$. So the column exists precisely when some $x\ge1$ in
the class has $z=(2ef-kb-xe)/f\ge1$, that is $xe+f+kb\le 2ef$. The left side increases with $x$, so
it suffices to test $x=x_k$, and then $x=x_k$ by Proposition~\ref{C-prop:closepaircolumns}(iv).
Writing $x_k=ke-m_kf$, the identity
\[
(ke-mf)e+f+k(f^2-e^2)-2ef \;=\; f\bigl(kf-(m+2)e+1\bigr)
\]
turns that inequality, after dividing by $f>0$, into $kf\le(m_k+2)e-1$, i.e.\ $g(k)\ge0$.

(ii) From $x_k=ke-m_kf$ and $x_{k+1}=(k+1)e-m_{k+1}f$ one gets
$x_{k+1}-x_k=e+(m_k-m_{k+1})f$. Both $x_k,x_{k+1}$ lie in $[1,f]$, so the left side lies in
$[1-f,f-1]$; with $0<e<f$ this forces $m_k\le m_{k+1}\le m_k+1$. Hence
$g(k+1)-g(k)=(m_{k+1}-m_k)e-f\le e-f<0$.

(iii) Since $1\le e\le f-1$, the least positive member of the class of $e$ is $e$ itself, so
$m_1=0$ and $g(1)=2e-1-f$. If $f\ge2e$ then $g(1)<0$, and $g$ is decreasing, so $g(k)<0$ for every
$k\ge1$. In general $g(k)\le g(1)-(k-1)(f-e)$, and $g(K)\ge0$ gives the stated bound on $K$.
\end{proof}


\begin{remark}[The wedge-filler hypothesis, and where it fails]\label{rem:alphamin}
Recorded in \cite{Obstructions}. Full discussion in \cite{Obstructions}.
\end{remark}

\begin{lemma}[A $c$-corner carries a side $a$-edge]\label{lem:ccornerside}\lab{PROVED}
The base corner angle is $\beta$, with flanks $\{a,c\}$. If a corner tile lays $c$ on the base it
lays $a$ on the side, so that side's $a$-count $P'$ is positive; by Lemma~\ref{C-lem:sidequant}
$P'=fp$, hence $P'\ge f$, and with $p e+R'=f$ and the $\gamma$-trap $R'\ge1$ one gets
$1\le p\le(f-1)/e$. So the side condition of Hypothesis~\ref{C-hyp:walls}, namely $p=0$, fails at
any $c$-corner: under that hypothesis both base ends are $a$-edges, exactly as
Theorem~\ref{C-thm:e1reduce} concludes at $e=1$. What remains for $e\ge2$ is to exclude
$1\le p\le(f-1)/e$, which is the analogue of the entire $e=1$ programme rather than a single
lemma; the boundary alone does not do it, since $a$ and $c$ are both available at the first base
position once the walls column is in force.
\end{lemma}

\begin{proposition}[The chain machinery is member-independent]\label{prop:unify}\lab{PROVED}
Every ingredient of the corner chain holds at all members $(e,f)$, not only at $e=1$:
\begin{enumerate}
\item[\textup{(i)}] the branch relations $3\alpha+2\beta=\pi$ and $\gamma=2\alpha+\beta$ define
the family, so the unique fills of $\beta$, $\alpha$ and $\alpha+\beta$ (\texttt{fill\_beta},
\texttt{fill\_alpha}, \texttt{fill\_alpha\_beta}) and the straight and full vertex
classifications are unchanged;
\item[\textup{(ii)}] $\cos\gamma=-e/2f<0$, so $2\gamma>\pi$ and every reflected partner dies at a
straight junction, at every member;
\item[\textup{(iii)}] $b\sin\gamma=c\sin\beta$, so an $a$-up tile spans its strip, at every member;
\item[\textup{(iv)}] for $1\le j<f$ the only decomposition of $j\,b$ is $j$ copies of $b$. Writing
$j=y+z+w$, the equation gives $e^2(z+w)+x\,ef=w f^2$, so $f\mid e^2(z+w)$ and hence
$f\mid z+w$ by coprimality; since $z+w\le j<f$ this forces $z=w=0$, so $y=j$ and $x=0$. Both
reductions are derived rather than assumed (\texttt{partition\_jb\_gen}). Verified for all
coprime $(e,f)$ with $e<12$, $f<20$.
\end{enumerate}
The single parameter that changes is the overshoot $c-b=e^2$, which is $1$ only at $e=1$. The
$e=1$ arguments that used $c-b=1$ as a unit therefore carry over with $e^2$ in its place, while
those that used its being the smallest positive integer do not.
\end{proposition}

\begin{lemma}[The $\alpha$-vertex gap, general]\label{lem:avgen}\lab{PROVED}
A cover piece of a horizontal $a$-edge has its foot at $(k+1)a$ from the row origin if it is
direct, and at $(k+3)a-e^3/f$ if it is mirrored, since $2\,cu_x-3a=-e^3/f$. The mirrored foot is a
lattice point only if $f^2\mid e^2$, and at row $1$, where feet land on the base and every
junction there is at an integer abscissa, it is a junction only if $f\mid e^3$; coprimality forces
$f=1$ in both cases (\texttt{alpha\_vertex\_gap\_gen}). So the mirrored placement is excluded at
row $1$ at every member with $f\ge2$, with $e^3$ playing the role that $1$ plays at $e=1$.
\end{lemma}

\begin{remark}[The two corners are not equally rigid]\label{rem:cornerasym}
At a $c$-corner the side junction carries $\gamma$ and the remaining $\alpha+\beta$ has the unique
fill $\{\alpha,\beta\}$, so the partner is forced immediately; this is what makes
Theorem~\ref{C-thm:n1} short. At an $a$-corner the side junction carries $\alpha$ and the remaining
$2\alpha+2\beta$ admits both $\{2\alpha,2\beta\}$ and $\{\beta,\gamma\}$, so the side branches at
once. The base end behaves the other way: there the $a$-corner leaves $\alpha+\beta$, whose unique
fill again forces the partner and confines the next base edge to $\{a,c\}$. The chain therefore
runs along the base from either corner, and what it cannot do is pass a base $b$-edge; this is the
mechanism behind Theorem~\ref{C-thm:e1cascade} at $e=1$ and it is available verbatim at every member.
\end{remark}

\begin{lemma}[The corner step, unconditional at every member]\label{lem:cornerstep}\lab{PROVED}
Let a base corner tile lay $a$ on the base. It presents $\gamma$ at the far end of that $a$-edge,
where the junction is straight, so the remainder is $\alpha+\beta$ with the unique fill
$\{\alpha,\beta\}$. The corner tile's $b$-edge joins a base vertex to a side vertex, so both of
its ends lie on the boundary and the tile across it lays a whole $b$; that tile presents $\alpha$
at the junction, since a straight junction carries at most one $\gamma$. Hence the $\alpha$-tile
is adjacent to the corner tile and the $\beta$-tile is adjacent to the base, and the second base
edge, lying among the $\beta$-tile's flanks, is an $a$ or a $c$.
\end{lemma}

\begin{remark}[Where the chain stops being free]\label{rem:chainobstruction}
Recorded in \cite{Obstructions}. Full discussion in \cite{Obstructions}.
\end{remark}

\begin{remark}[The thick program]\label{rem:thickprogram}
Recorded in \cite{Obstructions}. Full discussion in \cite{Obstructions}.
\end{remark}

\begin{remark}[First thick-member certificates]\label{rem:thickcerts}
The full $m=1$ search of the $(2,3)$ target ($N=23$, prime, $23\equiv11\pmod{12}$) is
\textsc{exhausted} at $19{,}677$ nodes: no tiling exists, and Hypothesis~\ref{C-hyp:walls} holds at
$(2,3)$. The refutation profile differs from the thin members: the semigroup-run prune carries
$57\%$ of the kills (against $8\%$ at $(1,5)$), consistent with Lemma~\ref{C-lem:monochotomy}
strengthening the run arithmetic, and the maximal depth is $5.7f$ against $4.8f$. \end{remark}

\begin{lemma}[The crossing kill, and solitude of branches]\label{lem:solitary}\lab{PROVED}
A transversal edge crossing a ladder's line within parameter $b+c$ of its start is fatal: the
covers cannot cross it, every T-junction forces the next cover edge, and
Lemma~\ref{C-lem:termination} forbids termination below $b+c$. A branch's own ceilings never
trigger this --- the $k$-th strip crossing is endpoint-tangent, sitting exactly $(k-1)c$ east of
the corresponding ceiling's end (\texttt{crossing\_tangency}) --- but a \emph{different} same-row
branch's ceilings do whenever it sits $1$ to $2f-1$ columns west: the offset $D f$ lands the
crossing inside a ceiling for every $-2f<D<0$, at parameter $b$ or $2b<b+c$. Hence any strip
junction of an $A$-branch with another $A$-branch $1$ to $2f-1$ columns to its west admits no
straddler, and its fills are mates: at most one $A$-branch per $2f$-column window per row-band
carries any straddler at all.
\end{lemma}

\begin{lemma}[The mirrored piece is charged]\label{lem:charge}\lab{PROVED: the junction arithmetic VERIFIED}
A cover piece of a horizontal $c$-edge is direct, with foot at the row lattice, or mirrored, with
foot displaced by $1/f$; the mirrored piece presents $\beta$ at its left junction, as does its
left neighbour, so that junction carries $n_\beta\ge2$ and its figure is $\{3\alpha,2\beta\}$
(\texttt{mirrored\_left\_junction}). By the census (Lemma~\ref{C-lem:census}) each such junction
forces an additional $\{\beta,3\gamma\}$ vertex: a single mirrored piece already gives $v_1\ge2$
(\texttt{escape\_charge}). The off-lattice escape available at rows $j\ge3$ is therefore paid for
globally, in climber vertices, of which the kernel-checked tilings carry between one and four.
\end{lemma}

\begin{lemma}[Climbers detect deviation]\label{lem:climberdetect}\lab{PROVED}
At an interior lattice point of the brick/mate structure exactly six tiles meet, contributing
$\beta,\gamma,\alpha$ from the row above and $\alpha,\gamma,\beta$ from the row below; the counts
are $(2,2,2)$ and the figure is $\{2\alpha,2\beta,2\gamma\}$. Hence a $\{\beta,3\gamma\}$ vertex
never occurs where the six incident tiles are the brick/mate six: by Lemma~\ref{C-lem:census} every
tiling contains such a vertex, and by Lemma~\ref{lem:charge} each mirrored cover piece forces a
further one, so each escape requires two points at which the tiling deviates from the forced
pattern. This is a detector, not an exclusion: the kernel-checked tilings do carry climbers at
lattice points, at places where their local structure is not brick/mate.
\end{lemma}

\begin{remark}[What reach buys, and what it does not]\label{rem:reachlaw}
Recorded in \cite{Obstructions}. Full discussion in \cite{Obstructions}.
\end{remark}

\begin{remark}[The residue at $f\ge6$]\label{rem:f6residue}
Recorded in \cite{Obstructions}. Full discussion in \cite{Obstructions}.
\end{remark}

\begin{remark}\label{rem:e1residual}
The residual class for $f\ge5$ is precisely: some side's first $a$-run behind an initial $c$-block
of length at least $2$. At $(1,4)$ the single such side word dies by the line-total kill of
Remark~\ref{C-rem:walls14}, whose three gap facts $b-a$, $b-e^2$, $2b-a\notin\langle a,b,c\rangle$
hold for every $f\ge3$; what does not yet generalize is the case analysis behind them when the
initial block exceeds two or the run splits. Closing this class for all $f$ is the last step of
Conjecture~\ref{C-conj:advance} at $e=1$.
\end{remark}

\section{Scope limits and closed routes from the main paper}\label{sec:mainscope}

\begin{remark}[Why the colouring method cannot close the rest]\label{OM-nogo}

The residual gap ($e\ge2$) is not merely unproved by \cite[Thm.~14]{BeesonIII}: it is
\emph{unprovable by that method}. Let $\eta=e^{i(\omega+\pi/2)}$, so that $f\eta^2+e\eta+f=0$ and
every edge direction of every such tiling lies in $\langle\eta,-1\rangle$. Any $T$-junction-proof
invariant of the shape $\sum_{\text{edges}}L\cdot\chi(\theta)$ with $\chi(-w)=-\chi(w)$ --- i.e.\
exactly the $\Phi$/colouring method, and exactly Beeson's --- is equivalent to the master identity
\[
m\bigl(-c\,z^4+a\,z^3+b\,z^2+a\,z-c\bigr)=z^3(f-ez)P(z)+(fz-e)Q(z),
\qquad P,Q\in\Z[z,z^{-1}],
\]
whose $z=\pm1$ specializations return exactly $M=m(f+e)$ and $-m(f-e)$. But that identity is
\emph{solvable for every} $m$, with the explicit witness $P_0=mf(z^{-1}-z^{-3})$,
$Q_0=m(ez^2-fz^3)$. Hence no invariant in this class can force $g\mid M$, for any $(e,f,m)$: the
$\Phi$-method --- the engine of both Beeson's argument and of our own $2\pi/3$ theorem --- provably
cannot decide the base-$\beta$ residue. Closing it requires genuinely non-linear input.

The natural non-linear candidate --- the full Conway--Lagarias tiling group, whose abelianization is
exactly this $\Phi$ class --- does not help either. For the base-$\beta$ tile $(2,3,4)$ we computed
that group over the finite quotients $\mathbb F_p[i]$ ($p\equiv3\pmod4$, $15$ a quadratic residue),
across $p\in\{43,59,67,71,103,127\}$: the class-$2$ (commutator) layer collapses completely onto the
tile relators, so the group reduces to its rank-$2$ abelianization $(M_\alpha,M_\beta)$, which is
admissible for every candidate. In particular the $m=1$ boundary word for $N=11$ is trivial in the
group --- indistinguishable from the $m=2,3,4$ words that \emph{do} tile --- while the controls hold
(a triangle not congruent to the tile is non-trivial, and the tileable $N=44,99,176$ words are
trivial). So the tiling group is not an obstruction here, extending the collapse we observed for the
$(8,7,13)$ tile (Remark~\ref{M-rem:tilinggroup}) to the base-$\beta$ tile itself
(\texttt{code/verify\_tiling\_group\_basebeta.py}). The non-linear input that closes the branch, if
any, is not homological.

Finally, the standard \emph{interior} tool fails here too, and for a reason that is an identity
rather than an accident. Beeson's $\Gamma_c$ argument --- the engine of his equilateral no-prime
theorem --- runs by showing $\Gamma_c$ is empty: if no relation $jc=ua+vb$ is witnessed with $j$
below the tile count, a $\Gamma_c$ link must emanate from the boundary, which is impossible since
links have in- and out-degree $1$ and cannot terminate on $\partial ABC$. For the base-$\beta$ tile
that hypothesis is false uniformly. Reducing $jc=ua+vb$ modulo $f$ twice gives $f\mid v$, then
$f\mid u-we$, and the complete solution set
\[
v=fw,\qquad u=we+fs,\qquad j=se+wf \qquad (w,s\in\Z_{\ge0}),
\]
generated by
\[
e\,c=f\,a \quad (j=e), \qquad\qquad f\,c=e\,a+f\,b \quad (j=f).
\]
The first is the identity $e f^2=f\cdot ef$. So a relation with $j=e$ always exists, and
$e<f\ll N=3f^2-e^2$; the $\Gamma_c$-emptiness contradiction therefore never triggers, for any
$(e,f)$. The classification and the witness are machine-checked in Lean
(\texttt{lean/GammaC.lean}: \texttt{gammac\_classification}, \texttt{gammac\_witness},
\texttt{gammac\_j\_lt\_N}; axiom-clean). Closing this branch needs control of a \emph{non-empty} $\Gamma_c$, which is exactly the
step left open in \cite{BeesonIso}.

The three failures above share one feature, and that feature settles what any future argument may
look like. Fix the tile $(2,3,4)$, which is the base-$\beta$ tile at $(e,f)=(1,2)$, and consider its
targets $(8m,8m,11m)$ with $N=11m^2$. These targets are pairwise similar and the tile is the same in
each. The case $m=1$ ($N=11$) admits no tiling, by exhaustive search; the cases $m=2$ and $m=3$ do
admit one, by the kernel-checked certificates of Remark~\ref{M-rem:thm14false}. Hence, if $\mathcal C$
is any criterion for non-tileability whose hypotheses depend only on the similarity class of the tile
and that of the target, taken separately, then $\mathcal C$ cannot exclude the base-$\beta$ target at
$m=1$: its hypotheses hold verbatim at $m=2$, where a tiling exists, so $\mathcal C$ would contradict
it. The statement is formalized and axiom-free
(\texttt{lean/ScaleRigidity.lean}: \texttt{no\_similarity\_invariant\_proof}); only the
\emph{existence} at $m=2$ is used, never the search result at $m=1$.

Angle arithmetic, the direction group $\langle\eta,-1\rangle$, the $\Phi$ class, the angular layer of
the Conway--Lagarias group, and any Dehn-type invariant are all criteria of this kind. So the three
obstructions recorded above are not three accidents but one, and the same applies to any further
invariant assembled from the same data. What survives is exactly the scale-sensitive information
carried by the edge-walk conditions, which do separate the cases: for the tile $(2,3,4)$ every
admissible walk on an equal side is free of $b$-edges when $m=1$, whereas at $m=2$ the walk
$(P,Q,R)=(3,2,1)$ is admissible and carries two of them (both machine-checked, ibid.). Any proof of
the base-$\beta$ no-go must take such input.
\end{remark}

\begin{remark}[A parity argument that fails, and why]\label{OM-parityfail}

It is tempting to argue as follows. Every tile has exactly one $b$-edge; by
Proposition~\ref{M-prop:unsplit} the edge $b$ is never a sum of two or more tile edges; so an interior
$b$-edge is matched by exactly one further $b$-edge, interior $b$-edges pair the tiles two at a time,
and $N\equiv Q_\partial\pmod2$ with $Q_\partial$ the number of boundary $b$-edges. We record that this
is \emph{false}, because the middle step does not follow.

Unsplittability forbids an edge from being \emph{exactly partitioned} by two or more tile edges. It
does not forbid a \emph{straddle}: a longer edge may contain a $b$-edge in its interior and overhang
it at one or both ends, and then the $b$-edge has no partner. Indeed the apex-mismatch configuration
of \cite{Companion} is built on precisely such a straddle. The genuine tilings confirm the
failure: in the $44$-tiling six $b$-edges belong to a single tile although \emph{no} $b$-edge lies on
the boundary, and the $99$-tiling has thirty-three matched pairs, eleven boundary $b$-edges and
twenty-two interior straddles. The correct relation is
\[
N=2\cdot\#\{\text{matched pairs}\}+\#\{\text{boundary}\}+\#\{\text{straddled}\},
\]
so the parity conclusion needs the straddle count to be even, which nothing supplies. (It happens to
be even, $6$ and $22$, in these two tilings, which is why the conclusion survives there.)

The same objection defeats the analogous arguments for the $a$- and $c$-edges. We state this
explicitly because the argument is natural, its conclusion is correct on every example we can test,
and it is nevertheless not a proof.
\end{remark}

\begin{remark}[the tiling group collapses as well]\label{OM-tilinggroup}

The strongest standard combinatorial obstruction is the Conway--Lagarias \emph{tiling group}
$T=\langle\chi_v\rangle/\langle$ collinear additivity;\ tile boundary words $=1$ in every orientation
and reflection $\rangle$; a tileable region $R$ has $\chi(\partial R)=1$ in $T$. Writing edge vectors
as complex numbers in $K=\Q(\sqrt3,i)$ (rotation by $\alpha$ is multiplication by
$\zeta_\alpha=(11+4\sqrt3\,i)/13$, by $\tfrac\pi3$ by $(1+\sqrt3\,i)/2$, reflection by conjugation)
and reducing to finite quotients $\mathbb F_{p^2}$ with $p\equiv11\pmod{12}$ --- so $\sqrt3\in\mathbb
F_p$ but $i\notin\mathbb F_p$ and the plane stays genuinely two-dimensional --- the tile relators,
\emph{together with} the commutators $[\text{class-1}(\text{tile}),\,\text{generator}]$ that normal
closure forces (a tile word's class-1 image is itself a relator, hence not central), fill the entire
commutator layer. So $T$ collapses to its rank-two abelianization $(M_\alpha,M_\beta)$: the class-2
quotient is trivial while the class-1 layer keeps rank $2$. As $N=105$ is admissible,
$\chi(\partial_{105})=1$, and the tiling group does not obstruct it. We verified this exactly for
$p\in\{23,47,59,71,83,107,131\}$ (\texttt{verify\_tiling\_group.py}), and --- extending to the free
nilpotent quotient through class $4$ and to a range of finite non-abelian groups (through $S_5$,
$\mathrm{GL}_2(\mathbb F_p)$, $\mathrm{PSL}(2,7)$ and the relevant Frobenius groups, which carry
genuine non-abelian data: sensitivity controls fire on $92$--$96\%$ of tile-consistent
representations) --- found \emph{no} tile-consistent quotient sending $\partial_{105}$ off the
identity, with the reptile controls $N=4,9$ trivial throughout and a $\Phi$-nonadmissible loop not.
The no-go of Remark~\ref{M-rem:conserved} therefore extends past the linear invariants to the full
tiling group: on this tile the tiling-group obstruction is exactly $\Phi$, so deciding realizability
of $N=105$ lies outside every standard tiling invariant.
\end{remark}

\begin{remark}[The scalene $2\pi/3$ shape $F_4$ first becomes possible at exactly $88$]\label{OM-F4at88}

The sweeps of Theorems~\ref{M-thm:frontier2}--\ref{M-thm:frontier4} test the isosceles and $F_1$ targets
among the $2\pi/3$ shapes and omit the scalene families $F_2,F_3,F_4$, on the ground that their
minimal counts exceed the range under consideration. That ground is exact, and it expires inside the
present band. Writing $N=k^2N_0$ with $N_0$ the closed-form ratio of
Section~\ref{M-sec:full}, the minima over primitive $120^\circ$-triples are
\[
\min N_0(F_2)=143 \ \text{ at } (a,b)=(3,5), \qquad
\min N_0(F_3)=264 \ \text{ at } (5,3), \qquad
\min N_0(F_4)=88 \ \text{ at } (3,5),
\]
so in $N\le 90$ these three families contribute exactly one instance, and it occurs at the endpoint
$N=88$: the tile $(3,5,7)$, with its $2\pi/3$ angle opposite $7$, on the scalene target
$(21,55,56)$, at $k=1$. The area ratio is exact, $660\sqrt3 = 88\cdot\tfrac{15\sqrt3}{2}$. It is the
first instance from any of $F_2,F_3,F_4$ to enter the frontier sweep --- scalene targets themselves
are not new, the sweep of Section~\ref{M-sec:full} having already searched the same tile $(3,5,7)$ on
$(15,21,24)$ at $N=24$ and $(7,8,13)$ on $(56,104,120)$ at $N=120$ --- and it is the reason $88$ is
not settled by the four exhaustions above. The instance is under search. The bound also confirms
that no scalene $2\pi/3$ shape can affect any $N<88$, so Theorem~\ref{M-thm:frontier4} and everything
below it are untouched.
\end{remark}

\begin{remark}[The third isosceles base-$\alpha$ instance]\label{OM-isoalpha84}

The isosceles base-$\alpha$ shape of \cite[Thm.~19]{BeesonIII} must be enumerated \emph{without} that
theorem's divisibility condition $g\mid M$, whose proof is unsound (Remark~\ref{M-rem:b3prime}):
imposing it reports the branch empty at $N=21$ and $N=70$, whose instances are real and needed the
engine. Solving the branch's defining cubic
$N/M^2=(1+s)(2-s^2)/\bigl((1-s)(2+s)^2\bigr)$ over the rationals in $0<s<1$, the solutions with
$N\le 90$ are exactly three:
\[
(N,M,s)=(21,5,\tfrac12),\qquad (70,8,\tfrac23),\qquad (84,10,\tfrac12),
\]
the first two being the instances of Theorem~\ref{M-thm:frontier2} and Remark~\ref{M-rem:seventy},
exhausted at $13\,659$ and $134\,631\,158$ nodes. The third is new and is the only one above $80$.
Since $s=e/f$ it forces $(e,f)=(1,2)$ and hence the tile $(2,3,4)$ --- the same tile as the
$44$-tiling --- on the isosceles $(24,24,42)$ with base angles $\alpha$, apex at $(21,3\sqrt{15})$.
So $84$ carries three instances, not the two the $(a+b)$-shape sweep produces: $(10,21,25)$ on
$(210,210,84)$, already exhausted at $78$ nodes, and $(110,21,121)$ on $(462,462,420)$, exhausted at
$4\,347\,701$ nodes, leaving only $(2,3,4)$ on $(24,24,42)$ under search. Sparse though this branch is --- three instances in
ninety values --- it is not exhausted by them, which is why the band cannot be closed by $N$-forms
alone.
\end{remark}

\begin{remark}[no invariant decides realizability]\label{OM-conserved}

The directional invariants are \emph{conserved} by the tiling process: each tile carries value
$\pm(c+a-b)$ for $f_\alpha$ (resp.\ $\pm(c+b-a)$ for $f_\beta$), so placing one tile changes
$\Phi_f$ of the unfilled region by $\pm V$ and the number of remaining tiles by one --- in step,
so that $\Phi_f/V$ stays integral and keeps the parity of the tile count throughout. Hence the
only conditions such an invariant imposes are the admissibility conditions of
Theorem~\ref{M-thm:spectrum}, and \emph{no} additive boundary invariant can distinguish an admissible
target from a partially tiled one; none can obstruct an admissible instance, and none prunes the
search below its combinatorial size. For the sharpest open instance $N=105$ we confirmed this
computationally: searching all directional weights of period $\le4$ in the $\alpha$-winding
(reflections included), the only invariants are $f_\alpha$ ($M_\alpha=5$) and $f_\beta$
($M_\beta=21$), both admissible; and the vertex-type balance forced by the angle relations
$a\alpha+b\beta+c\gamma\in\{\pi,2\pi\}$ admits $42{,}282$ nonnegative integer solutions. Deciding
realizability of an admissible instance therefore lies beyond every linear or coloring invariant;
it is a genuinely combinatorial question, which is why the search of Theorem~\ref{M-thm:decidable}
appears to be essential.
\end{remark}

\begin{remark}[The spectrum is member-dependent]\label{OM-memberdep}

Theorem~\ref{M-thm:realize12} does not extend along the base-$\beta$ family. For $(e,f)=(1,3)$ --- tile
$(3,8,9)$, $N_1=26$ --- the scale-$2$ target admits no tiling: $N=104$ is excluded by exhaustive
search ($43\,541\,916$ nodes). The two ingredients of the induction therefore behave differently
along the family: the unit parallelogram is constructible for every member with $e=1$ (companion
note, Theorem on the unit parallelogram; verified for $f=2,\dots,12$), while the seeds $\Delta_2$,
$\Delta_3$ that the collar needs are not. Every attempt to build the seed at $f=3$ fails for the same
reason, namely that there is nothing to build: the cevian piece, the symmetric middle piece, and the
maximal interiors of both natural lattices all return \textsc{no tiling}, and a seeding by unit
parallelograms is impossible by counting. Consistently with this, $(1,2)$ is the unique member of its
family with $\beta<\pi/3$. What the spectrum of $(1,3)$ is remains open: since $\Delta_4$ would
require $\Delta_1+\Delta_3$ and $\Delta_5$ would require $\Delta_2+\Delta_3$, both excluded, the set
reachable from $\Delta_3$ alone is exactly the multiples of $3$.
\end{remark}

\begin{remark}[Why the forms, and what the reformulation buys]\label{OM-forms}

Every tile count in the classification is a value of a binary quadratic form: $e^2+f^2$
(discriminant $-4$) on the commensurable branch, the Eisenstein norm $a^2+ab+b^2$ (discriminant
$-3$) for a $2\pi/3$ tile, $3f^2-e^2$ (discriminant $12$) and $2f^2-e^2$ (discriminant $8$) on the
$3\alpha+2\beta$ branch. The classical exclusion of primes $\equiv3\pmod4$ on the commensurable
branch is exactly Fermat's two-squares theorem for the first form; Theorem~\ref{M-thm:mod12} is its
analogue for the third. Two consequences. First, the exception in Theorem~\ref{M-thm:main} is not a
Diophantine side condition but a single arithmetic progression, and by Dirichlet it carries half of
the primes $\equiv3\pmod4$: the theorem is unconditional on the other half. Second,
$19\equiv7\pmod{12}$, so the value in Erd\H os's question lies in the unconditional class.
The forward implication is machine-checked (\texttt{basebeta\_prime\_mod\_twelve},
\texttt{BaseBetaMod12.lean}); it is the direction the corollary uses. The converse was checked for
all $109$ primes $\equiv11\pmod{12}$ below $3000$ with no exception.
\end{remark}

\begin{remark}[What $83$ rests on]\label{OM-eightythree}

Among the values of the band, $83$ is of a different kind from the rest, and the distinction is
sharper than mere corroboration. Every prime $\equiv11\pmod{12}$ below it --- $11,23,47,59,71$ ---
was settled by an independent certified exhaustive search of its base-$\beta$ instance, so those
exclusions stand on the engine alone and need no branch theorem. For $83$ the single surviving
instance is the tile $(1722,83,1764)$ on $(3486,3486,3403)$, the base-$\beta$ member $(e,f)=(5,6)$,
and that search has not terminated. Its exclusion therefore rests entirely on
Theorem~\ref{M-thm:fullprime}, which is conditional on the companion's complete-corner-wall
hypothesis (Remark~\ref{M-rem:conditional}). \emph{So $83$ is not excluded unconditionally by this
paper}, and it is the first value for which that is true: everything below $81$ is unconditional,
either by an $N$-form, by Theorem~\ref{M-thm:main} for primes $\equiv7\pmod{12}$, or by an exhaustion.
We state this explicitly because the boundary between the two kinds of result is easy to lose in a
table, and because the same hypothesis governs every $N\equiv11\pmod{12}$ above $83$ as well.
\end{remark}

\begin{remark}[What this does and does not remove]\label{OM-productscope}

Proposition~\ref{M-prop:product} is a statement about the $2\pi/3$ branch only. It says nothing about
the base-$\beta$ branch of $3\alpha+2\beta=\pi$, which is a different branch with a different tile
and is where the one genuine gap of Theorem~\ref{M-thm:main} lives (Remark~\ref{M-rem:isobeta}). What it
does buy is a narrowing of the rationality dependence: Corollary~\ref{M-cor:similar} and
Proposition~\ref{M-prop:F1free} discharge three of the eleven shapes --- tile-similar, $F_1$, $F_1'$
--- with no appeal to \cite{BeesonZhang} or \cite[Thm.~12.5]{BeesonIso}, whereas the remaining eight
still use an integer-sided tile. The obstruction to doing the same for $F_2,F_3,F_4$ is visible in
the table: there $\kappa\ne\pm1$, so the product pins only the single ratio $a/b$, and the law of
cosines then leaves $c/b$ a quadratic irrationality that no further invariant condition controls.
Establishing rationality on a general target remains exactly the content of the cited theorem.
\end{remark}

\begin{remark}[Scope]\label{OM-ratfreescope}

The hypothesis is explicit and strictly weaker than the imported theorem: excluding $(\star)$ is a
corner-and-walk exclusion of the kind carried out for the base-$\beta$ target, where the
$\gamma$-trap shows every side carries a $c$-edge. It is not established here for these six shapes.
The remaining two shapes are unaffected: on $F_2$ the same elimination gives a quadratic in $a/b$
with discriminant $9\kappa^2-30\kappa+9$, and on the equilateral target $\kappa=3$ is constant, so
both continue to use \cite{BeesonZhang}. The arithmetic is machine-checked in
\texttt{lean/RationalityFree.lean}. That the condition has content is visible from the fact that
$a/b\in\Q$ alone does not give $c/b\in\Q$: the latter asks that $(a/b)^2+(a/b)+1$ be a rational
square, which holds for $a/b=3/5$ and $7/8$ but not for $a/b=1$ or $2$; the tiles for which it holds
are the integer triangles with a $2\pi/3$ angle.
\end{remark}

\begin{remark}[Scope]\label{OM-rationality}

The proposition converts the rationality input from a hypothesis about the tile into a hypothesis
about the invariant, and the latter is what the method produces natively: wherever the boundary
functional of Section~\ref{M-sec:invariant} is defined and integral, rationality is automatic. It does
not by itself discharge the citation on the \emph{general} target, because establishing integrality
of all three quantities there is exactly the content of the cited theorem. What it does discharge is
the appearance of a rationality assumption \emph{inside} the invariant argument, where the three
integralities are already in hand; and it shows the two facts are equivalent rather than
independent, which is not how the literature presents them.
\end{remark}

\begin{remark}[Scope]\label{OM-ladder}

Two points. First, the exclusion of $11$ together with the realizability of $11\cdot2^2$ and
$11\cdot3^2$ shows that a base-$\beta$ family can be empty at scale $1$ and non-empty above it; the
obstruction at $m=1$ is genuinely a small-scale phenomenon and not a property of the family. Second,
before this the smallest tiling known in any incommensurable branch had $1215$ tiles; the corollary
supplies infinitely many, of which $44$ is the smallest. Whether $S(e,f)\ne\emptyset$ for every
$(e,f)$ --- equivalently, whether every base-$\beta$ family tiles at \emph{some} scale --- is open;
the first untested instances are the scale-$2$ targets of the families $(2,3)$, $(3,4)$ and $(1,4)$,
with $N=92$, $156$ and $188$.
\end{remark}

\begin{remark}[The golden ratio]\label{OM-golden}

Two consequences. First, Proposition~\ref{M-prop:unsplit} is exactly the hypothesis the corner
parallelogram needs, so Proposition~\ref{M-prop:cornerpara} holds with \emph{no} proviso. Second, $e=1$ --- which
both genuine tilings realize --- is the only case in which an edge splits at all, while the open
$e\ge2$ regime is rigid. The comparison of $a$ with $b$ must be made with care:
$b>a\iff f^2-ef-e^2>0\iff f/e>\varphi$, the golden ratio. The thick regime straddles that threshold:
$N=66$ $(3,5)$ and $N=131$ $(4,7)$ are super-golden, while $N=23$ $(2,3)$, $N=59$ $(4,5)$,
$N=83$ $(5,6)$ and $N=167$ $(5,8)$ have $b<a$. Proposition~\ref{M-prop:unsplit} above assumes no
inequality between $a$ and $b$.
\end{remark}

\begin{remark}[What the alignment pins, uniformly]\label{OM-align}

For every open thick member the mismatch ray is now completely determined out to length $f\cdot b$: its
far wall is $b^{f}$, terminating transverse to the opposite side, and $V=f^2$ falls strictly inside the
$\lceil f^2/b\rceil$-th $b$-edge. Explicitly (all exact): $N=23$ far side $b^3=5^3$, $V=9\in(5,10)$;
$N=59$ $b^5=9^5$, $V=25\in(18,27)$; $N=66$ $b^5=16^5$, $V=25\in(16,32)$; $N=83$ $b^6=11^6$,
$V=36\in(33,44)$. This pins an extended \emph{interior} structure of a
hypothetical thick tiling, uniformly in $(e,f)$; it serves as a root prune for the apex-down search; the forced $b^f$ wall must land on the base at
coordinate $e\cdot f^2$, whose $b$-count is bounded (\texttt{base\_b\_bound}). It does not, on its
own, close any member.
\end{remark}

\begin{remark}[Scope]\label{OM-incomm}

The irrationality hypothesis cannot be dropped: the commensurable tile
$(\frac\pi6,\frac\pi6,\frac{2\pi}3)$ tiles an equilateral (hence isosceles) triangle with $N=3$, by
the fan from the centre; that tile lives in the commensurable branch of the classification, where
$N=3n^2$ is allowed. The equilateral target for an \emph{incommensurable} $2\pi/3$ tile is excluded
for primes by Beeson~\cite{BeesonEq}, and is covered in the assembly of Section~\ref{M-sec:assembly};
it is not part of this theorem.
\end{remark}

\begin{remark}[Scope]\label{OM-gammascope}

The hypothesis that the corner angles are $\beta,\beta,3\alpha$ is essential, and this is why
\cite[Lem.~6]{BeesonIII} is false as quoted (Remark~\ref{M-rem:lem6false}): on the tile-similar target
a corner \emph{is} a $\gamma$, the chain terminates harmlessly, and the counterexample there has two
sides with no $c$-edge. On a base-$\beta$ target no corner is a $\gamma$, and the argument runs.
\end{remark}

\begin{remark}[The field of computation]\label{OM-fieldshort}

For the tile $(ef,f^2-e^2,f^2)$ one has $\sin\alpha=e\sqrt D/(2f^2)$ with $D=4f^2-e^2=(2f-e)(2f+e)$,
so a base-$\beta$ tiling's coordinates lie in $\Q(\sqrt D)$. This is the field the exact search
computes in; for $(e,f)=(1,2)$ it is $\Q(\sqrt{15})$, the ring of all three tiling certificates.
\end{remark}

\begin{remark}[Validation on the $44$-tiling]\label{OM-rung2}

On the genuine $44$-tiling ($f^2=4>2e^2=2$, so $t=0$) the piercer spans exactly $[3,6]$, starting at
the T-junction, with corner $\alpha$ there to numerical precision, as forced; the arithmetic core is
machine-checked (\texttt{pre\_pierce\_dichotomy}). The rung also has a computational consequence: the
apex neighbourhood is rigid (three forced $\alpha$-tiles, a forced mismatch ray, a forced T-junction,
a forced piercer start), while the base admits several walks. Since the corner-anchored search
anchors at the lowest vertex, orienting the target apex-down (a reflection, under which verdicts are
preserved) places the rigid chain at the root of the search tree: on $N=23$ the apex-down search
exhausts in $3{,}954$ nodes against $19{,}677$, whereas on the thin member $N=47$, where the base is
the rigid side, the apex-down orientation is slower. The thick searches $N=59,66$ are run apex-down.
\end{remark}

\section{The fan-collapse prune}\label{sec:fanprune}

The searches that certify individual exclusions are exhaustive depth-first placements, and their
cost is dominated by one structure. At a convex corner of the uncovered region the covering tiles
must each have a \emph{vertex} there: a tile containing the corner in its interior leaves the
region, and one carrying it interior to an edge contributes $\pi$, which exceeds the corner. If the
corner is smaller than the tile's second-smallest angle, no angle but the smallest fits, and the
corner is completable only if it is an exact integer multiple of that angle.

\begin{proposition}[Soundness of the prune]\label{prop:fanprune}\lab{PROVED}
Let a convex corner of the uncovered region have angle $W$, and let the tile's angles be
$\alpha<\beta\le\gamma$. If $W<\beta$ then every tile angle at the corner equals $\alpha$, so
$W=k\alpha$ for $k$ the number of tiles there; consequently a corner with $W<\beta$ that is not an
integer multiple of $\alpha$ admits no completion, and the subtree below it is empty.
\end{proposition}

Machine-checked as \texttt{FanPruneSound.corner\_all\_alpha},
\texttt{corner\_is\_multiple} and \texttt{fan\_prune\_sound}. The engine tests the multiple
condition with $\cos k\alpha$ generated by the Chebyshev recurrence in exact rational arithmetic,
$\cos\alpha$ being rational in the side lengths. The prune removes no tiling, and the certified
$44$-tiling is still found with it enabled.

Its effect is to cut, at the root, subtrees the unpruned search explores through $2^{k}$ chirality
branches that all die for the same reason. On the $(4,2)$ orbit the node count falls from
$15\,953$ to $109$ at $f=8$, a factor of $146$; the orbit $(3,2)$ at $f=8$, which
had run over three million nodes without terminating, exhausts at $29\,733$. The growth in $f$
remains exponential --- measured at about $1.61$ per unit $f$ against $2.97$ without the prune, so
the base is roughly halved rather than the complexity class changed --- and the improvement is
uneven across orbits: at $f=12$ most exhaust within a few thousand nodes while $(6,5)$ does not
finish. What the prune bought is the exclusion of $N=191$; it does not by itself reach $N=431$.

Two further search improvements belong here, both sound and both measured. The corner test above
generalises: a convex corner is fillable exactly when its angle is $X\alpha+Y\beta$ for integers
$X,Y\ge0$, since the tiles at it contribute $x\alpha+y\beta+z\gamma=(x{+}2z)\alpha+(y{+}z)\beta$;
all of $\cos\alpha,\sin\alpha,\cos\beta,\sin\beta$ lie in $\mathbb Q(\sqrt D)$ for this family, so
the table of admissible corner cosines is computed exactly once per instance. And the vertex
ordering may be changed to fill the most constrained vertex first, which is sound for an
exhaustive search because every tiling fills every vertex, so the choice affects only the node
count. The second is dramatic in one direction and dangerous in the other, and the asymmetry is worth
recording. On the free-corner instances it is a large gain: $N=23$'s falls from $19\,677$ nodes to
$387$, $N=11$'s from $135$ to $42$, and the certified $44$-tiling is found in $155\,123$ rather
than $858\,163$, every verdict unchanged. On the word-constrained instances it is a severe loss:
the $(4,2)$ orbit at $f=12$ exhausts in $565$ nodes under the default ordering and fails to
terminate at all under the most-constrained one. The ordering is sound either way --- completeness
cannot depend on which vertex is filled first --- but it must be chosen per instance family, and
for the words identified below the default is the right choice. Neither improvement reaches those
words. A third does, and it is the one that matters. The fillability criterion above was being
applied only at \emph{acute convex} corners, because the corner routine skipped every non-convex
vertex and the criterion sat behind an acute-only guard. Yet the same equation binds at obtuse
convex corners, and at reflex corners as well, where the interior angle is $2\pi-\theta$ and the
cosine is that of $\theta$, so the test is made against a table for the second half turn. A first
implementation was unsound --- its table was built with a termination test on $\sin=0$, which for
irrational $\alpha/\pi$ never occurs, so angles beyond $2\pi$ entered the tables --- and the
positive control caught it at once, the certified $44$-tiling being reported as absent. With the
table bounded correctly by the sign of the sine, the certified tiling is found again and every
prior verdict and node count is reproduced, while the words that had resisted collapse: the two
$f=12$ words that never terminated exhaust at $506$ and $807$ nodes, the whole of $f=12$ in under
$150$ seconds, and $f=14$ likewise. That is what settles $N=431$ and $N=587$.

The criterion behind it is machine-checked. Every contribution at a corner --- a vertex showing
$\alpha$, $\beta$ or $\gamma=2\alpha+\beta$, or an edge through it contributing
$\pi=3\alpha+2\beta$, which can occur only at a reflex corner --- is of the form $X\alpha+Y\beta$
with $X,Y\ge0$, so a filled corner's angle is too, and a corner admitting no such representation
admits no fill (\texttt{FanPruneSound.corner\_fill\_form},
\texttt{corner\_unfillable}). The reflex case is matched against a separate table because there the
interior angle is $2\pi-\theta$ and $\cos(2\pi-\theta)=\cos\theta$, so the cosine alone does not
distinguish the two; the sign of the cross product does.

The residue is sharply located. Among the $f=12$ orbits, every word whose $c$ and $b$ letters are
not adjacent in the order $c$ then $b$ finishes within a few thousand nodes, while the single word
with $c$ immediately followed by $b$ does not finish at all; at $f=8$ the seven orbits contain no
such word and all finish in at most $743$ nodes, whereas the one $f=8$ word of that shape takes
forty times as long. The reason is structural: when $c$ and $b$ are adjacent the base splits into
two $a$-runs and one block of length $c+b=2f^{2}-1$, two thirds of the base, holding a single
internal junction --- and the propagation that pins the rest, namely that an $a$-edge's ends carry
$\beta$ and $\gamma$ while two $\gamma$s cannot meet, has nothing to act on inside it. The
$c\mid b$ junction is therefore the one junction type the orientation argument cannot reach, and
its analysis is what separates the present exclusions from $N=431$.

\section{A self-reducing family}\label{sec:march}

Completing the corner criterion did more than settle two values: it made the refutation trees small
enough to read, and reading them exposes a structure.

\begin{proposition}[The member-uniformity of the tests]\label{prop:testsuniform}\lab{VERIFIED}
Every test the search performs gives the same answer at every member. A corner of angle
$X\alpha+Y\beta$ is fillable exactly when $x+2z=X$, $y+z=Y$ has a nonnegative solution, a condition
naming no $f$ (\texttt{MemberUniform.fill\_iff\_counts}), so fillability transfers between members
in both directions (\texttt{fillable\_iff\_member}, \texttt{unfillable\_iff\_member}); and the
residue-$1$ test answers no at every member (\texttt{gap\_one\_uniform}), as does the
$|a-b|$ test (Lemma~\ref{lem:stubgap}). Hence no test the search performs can distinguish two
members (\texttt{tests\_uniform}).
\end{proposition}

\begin{proposition}[The adjacent family is independent of the member]\label{prop:findep}\lab{HEURISTIC: measured; the mechanism VERIFIED}
For a base word in which $c$ is immediately followed by $b$ --- that is, $cp=bp-1$ --- the
refutation is the same at every $f$: identical node count, identical depth, and identical values of
every prune counter. At $f=12,14,16,18,20,22$ the word $(6,5)$ gives $506$ nodes with depth $13$ in
every case, and $(5,4)$, $(7,6)$, $(8,7)$, $(9,8)$, $(10,9)$ give $299$, $807$, $1315$, $2124$,
$3441$ likewise. By Proposition~\ref{prop:testsuniform} no test the search performs can tell the
members apart, which is why; what remains measured rather than proved is that the search performs
only those tests.
\end{proposition}

\begin{proposition}[Its size is Fibonacci in $bp$]\label{prop:fib}\lab{HEURISTIC: measured; the recurrence's arithmetic VERIFIED}
Those sizes satisfy $a(n)=a(n-1)+a(n-2)+2$, equivalently $a(n)+2$ is a pure Fibonacci sequence
(\texttt{MarchRecurrence.shift\_to\_fibonacci}), with growth $\varphi$ --- the measured ratio being
$1.6181$ per unit $bp$. The recurrence holds at $bp=7,\dots,15$ against measured values, three of
them recorded earlier for an unrelated purpose, and its two untested values, $9010$ and $23591$,
were predicted before measurement and confirmed exactly.
\end{proposition}

The mechanism is visible in the traces. The forced spine is an $a$-tile, a second $a$-tile, the
$\alpha$-filler joining their two apexes, and a third $a$-tile --- the march along the $a$-run --- and
the branch is that filler's chirality, the two children differing by one unit in apex height. The
subtree sizes identify what the two choices do: one advances the march by a single position, giving
$a(n-1)$, the other by two, giving $a(n-2)$. It is a monomer--dimer march, which is why the count is
Fibonacci, and the $+2$ counts the spine nodes consumed at each step. The independence of $f$
follows: the march's length is set by $bp$, so the $a$-run beyond the $b$-letter never enters.

The decomposition is exact in the trees. $(6,5)$'s $506$ nodes split, after a four-node spine, into
subtrees of $205$ and $296$, where $205$ is the measured refutation of $(4,3)$ appearing verbatim
and $296$ is that of $(5,4)$ less its three shared spine nodes; $(5,4)$'s $299$ split as $92$ and
$202$ in the same way.

\begin{remark}[What would close the family]\label{rem:marchobl}\lab{OPEN}
An inductive proof needs the base cases together with one step at each $a\mid a$ junction of the
run. Three parts of that have now been discharged, and naming what is left is the point of this
remark.

\emph{Discharged.} The base cases are not the measured counts $205$ and $109$: the induction
bottoms out at runways of length $0$ and $1$, whose terminal stub is killed by residue $1$ being a
gap of $\langle f,f^2-1,f^2\rangle$ (\texttt{MarchStep.offset\_terminal\_dies}, with a non-vacuity
witness). The overshoot is at most one position (\texttt{.overshoot\_le\_one}). And the step's
\emph{vertex} half is now a theorem about a real dissection rather than a reading from the traces:
at a point of the frontier that is not a target vertex and at which some tile presents $\gamma$, the
figure is forced to $\{\alpha,\beta,\gamma\}$ --- one tile each, none straight
(\texttt{VertexFigureReal.gamma\_boundary\_figure\_real}) --- so the uncovered opening is $\alpha$,
and a tile filling it with its $\alpha$-corner there has its two edges on the wedge's two boundary
rays, in exactly two ways (\texttt{JunctionWedge.march\_junction\_real}). Earlier this consumed the
multiplicity equations as hypotheses; they are now derived from the boundary figure classification.

\emph{What is left, sharpened.} The three remaining statements are about the \emph{run} rather
than about a vertex, and they do not stand on the same footing.

\emph{(i) The march's steps land on vertices of the above kind.} This is now reduced to
Proposition~\ref{C-prop:orientmono} and needs nothing further. At a straight-edge point the figure
is $\{3\alpha,2\beta\}$ or $\{\alpha,\beta,\gamma\}$ and there is no third case
(\texttt{MarchRun.junction\_dichotomy}), so a junction either carries no $\gamma$ at all or carries
one and is a march junction (\texttt{.gamma\_of\_not\_exceptional}). Orientation monotonicity says at
most one junction of an $a$-run is of the first kind --- it is the $\mathrm{BG}\to\mathrm{GB}$
transition --- so all but at most one junction of the run is a march junction
(\texttt{.all\_but\_one\_is\_march\_junction}). The residual dependence is on
Proposition~\ref{C-prop:orientmono}, whose own blocker is bridge~(c).

\emph{(ii) and (iii) are blocked at the level of statement, not of proof.} The wedge's two rays
carry $b$ and $c$: at the $\gamma$-corner the flanking sides are $a$ and $b$, at the $\beta$-corner
$a$ and $c$ (\texttt{TilePlacement.incident\_sides}), and the filler's $\alpha$-adjacent sides are
$b$ and $c$. So the two chiralities are forced to be the \emph{flush} one, $b$ on the $b$-ray and
$c$ on the $c$-ray, and the \emph{crossed} one, leaving a discrepancy $c-b=1$ on each ray --- which
is why the two children differ by one unit in apex height. What does not follow is the displacement
\emph{along the run}: advance by one position versus two.

A coordinate model of the run now exists and settles part of (ii). Put the $a$-edge on the
$x$-axis from $(t,0)$ to $(t+f,0)$; the apex is the corner opposite $a$. At a $\gamma$-corner the
flanking sides are $a,b$ and at a $\beta$-corner $a,c$, so the apex lies at distance $b$ from the
$\gamma$-end and $c$ from the $\beta$-end, giving two orientations with horizontal offsets
\[
d_{\mathrm{GB}}=\frac{1-f^2}{2f},\qquad d_{\mathrm{BG}}=\frac{3f^2-1}{2f}=c\cos\beta .
\]
The second is \texttt{Rigidity}'s $x_w$ and matches the closed form $\cos\beta=(3f^2-1)/(2f^3)$ of
Lemma~\ref{C-lem:anglethreshold}. Three facts follow, all machine-checked
(\texttt{MarchCoords}):
\begin{enumerate}
\item[\textup{(a)}] Both orientations give the \emph{same} apex height
(\texttt{apex\_height\_common}); this is \texttt{FillerGeneral.filler\_b\_general} in disguise, and
it is why the branch is a chirality rather than a change of level.
\item[\textup{(b)}] $d_{\mathrm{GB}}+d_{\mathrm{BG}}=f$ (\texttt{offsets\_complementary}), so the
filler's $\alpha$-corner lands exactly on the junction and its two sides from there reach the two
apexes at distances $b$ and $c$ (\texttt{junction\_to\_apex\_bg}, \texttt{\_gb}). The chirality is
the $a$-tiles' orientation, which swaps which apex is at $b$ and which at $c$
(\texttt{chirality\_swaps\_sides}): the two placements of \texttt{MarchStep.two\_placements},
realised in coordinates as one triangle and its reflection.
\item[\textup{(c)}] Consecutive same-orientation apexes are exactly $a=f$ apart, so the segment
joining them is a tile edge; at an orientation change the spacing is $(3f^2-1)/f$, which is never a
tile side, since $3f^2-1=f\cdot s$ for $s\in\{f,f^2-1,f^2\}$ forces $2f^2=1$ or $f\mid 1$
(\texttt{spacing\_same}, \texttt{spacing\_gb\_bg}, \texttt{transition\_not\_a\_side}). So the
march's horizontal filler exists only where the orientation is constant. This is the local,
coordinate-side counterpart of Proposition~\ref{C-prop:orientmono}'s conclusion.
\item[\textup{(d)}] The filler is \emph{forced}, not chosen. Its three vertices are the junction and
the two apexes, so the two $a$-tiles determine it; and the triangle on those points has squared
sides $\{a^2,b^2,c^2\}$ as a multiset in either orientation, so it is congruent to the tile and the
configuration is admissible (\texttt{filler\_forced}). The branch is therefore the $a$-tiles' common
orientation and nothing else.
\end{enumerate}

What remains unproved in (ii) is the \emph{displacement along the run}: that the two chiralities
advance the march by exactly one and by exactly two positions. The configuration is now pinned in
coordinates, but the advance is a statement about which $a$-tile the next spine step begins from,
and that requires the forcing of the configuration along the whole run, not one junction. (iii) is
untouched.

One caution, recorded because it cost a session. Stating (ii) or (iii) by quantifying over an
abstract advance function is not merely unproved but false: ``$\mathrm{adv}(\text{flush})=1$ and
$\mathrm{adv}(\text{crossed})=2$ for every $\mathrm{adv}$'' is refuted by $\mathrm{adv}\equiv 0$, and
``$R(bp-1)\Rightarrow R(bp)$ for every $R$'' by $R(n)=[n\le 1]$ at $bp=2$; both refutations are
machine-checked. The advance is what the coordinates \emph{define}, so it must be read off the model
above, never quantified over.
\end{remark}

We record one further observation without explanation. The complementary family $cp=2$ is nearly
constant in $bp$ at fixed $f$ and grows in $f$ at $1.6164$, tending to $\varphi$; so both families
are golden-ratio, one in $bp$ and one in $f$. That points to the same march underlying both, with
only the parameter setting its length differing, but we could not identify which run plays that role
in the second family: the run before the $c$ has length one, and the run after the $b$ varies with
$bp$ while the measurements do not.

\section{Four macroscopic routes, closed}\label{sec:macro}

The routes recorded elsewhere in this paper are local. This section closes four \emph{global}
ones --- algebraic, spectral, topological and homological --- each proposed as a way to force the
tile count composite as $N\to\infty$. All four fail, and the reasons are worth having.

\begin{proposition}[The area ratio cannot exclude primes]\label{prop:norm}\lab{PROVED}
For the base-$\beta$ branch the area ratio is $N=m^{2}(3f^{2}-e^{2})$, and
$3f^{2}-e^{2}=-\mathrm{N}_{\mathbb Q(\sqrt3)/\mathbb Q}(e+f\sqrt3)$ is a norm form. Norms of prime
elements are prime, so no Galois or Diophantine obstruction in the ratio can force $N$ composite;
and indeed $144$ primes of this form occur with $\gcd(e,f)=1$, $e<f<40$. What the algebra supplies
is the \emph{identification} of the candidates --- a prime value forces $N\equiv11\pmod{12}$
(Theorem~\ref{M-thm:mod12}) with a unique representation --- not their exclusion.
\end{proposition}

\begin{proposition}[The global angle--Euler system admits prime solutions]\label{prop:globalsys}\lab{PROVED}
Independence of $\alpha,\beta$ over $\mathbb Q$ permits the vertex figures $(3,2,0),(1,1,1)$ at a
boundary non-corner point, $(6,4,0),(4,3,1),(2,2,2),(0,1,3)$ at an interior point, and the unique
corner fills. Each tile contributes one corner of each type, so each type sums to $N$; with the
Euler relation this yields
\[
  N \;=\; 2I + B + 1 ,
\]
and the Euler count and the angle count give \emph{the same} relation, no more. Solving the full
system over the nonnegative integers gives $17$ solutions at $N=11$, $172$ at $N=23$ and $1968$ at
$N=47$. Since $N=11$ is excluded by certified exhaustion, the global system is strictly weaker than
the truth, and its only congruence --- $N$ odd iff $B$ even --- holds for every prime.
\end{proposition}

\begin{remark}[Spectral invariants of the dual graph]\label{rem:spectral}\lab{PROVED}
The dual graph of our certified $44$-tiling, built from its exact coordinates, has $44$ vertices,
$45$ shared full edges, degree distribution $\{1^{11},2^{20},3^{13}\}$ and four components of
orders $4,8,16,16$ --- not isomorphic, so a decomposition into isomorphic subgraphs does not occur
in the one realized example available. More decisively, a spectral criterion is a function of the
dual graph, whose admissible constraints are exactly the vertex-figure counts of
Proposition~\ref{prop:globalsys}; those admit prime solutions, so no such criterion excludes a
prime order.
\end{remark}

\begin{remark}[No higher nilpotent layer]\label{rem:nilptower}\lab{PROVED}
The lower central series satisfies $\gamma_{k+1}=[G,\gamma_k]$. The tiling group's class-2 layer
collapses (Remark~\ref{OM-tilinggroup}), i.e.\ $\gamma_2=\gamma_3$; then
$\gamma_4=[G,\gamma_3]=[G,\gamma_2]=\gamma_3$ and inductively $\gamma_k=\gamma_2$ for all $k\ge2$.
Every nilpotent quotient therefore equals the abelianization: the obstruction is not that class $2$
was the layer chosen, but that the tower has no further layers.
\end{remark}

The four share a cause. Each is a counting invariant --- of areas, of spectra, of vertices, of
commutators --- and \S\ref{sec:onewall} records that counting cannot answer the crossing question.
Proposition~\ref{prop:globalsys} is the sharpest form of that statement we have: the complete
global system is satisfied, in seventeen ways, by a value proved untileable.

\section{The crossing question is the single obstruction}\label{sec:onewall}

The preceding sections record failures one route at a time. This section states the consolidated
fact they add up to, which is sharper than any of them separately: every route to the base-$\beta$
primes known to this programme reduces to one question, and every invariant class available to the
field is provably unable to answer it.

\begin{definition}[The crossing question]\label{def:crossQ}
Let $(e,f)$ be coprime with $1\le e<f$, let $\mathcal T$ be a tiling of the base-$\beta$ target at
$m=1$, and let $S$ be a maximal run of forced tiles whose far sides lie along a segment $L$ of a
line (a corner chain along a wall, or a column footprint along a level line), with $V$ the endpoint
of $L$ beyond the run's last forced tile. The question: \emph{must the next tile along $L$ lay a
whole edge on $L$ without overrunning $V$, or can it cross $L$ at $V$?}
\end{definition}

\begin{proposition}[Three reductions]\label{prop:threecostumes}\lab{PROVED}
\emph{(i)} At $e=1$ the entire subfamily --- every prime $3f^2-1$ --- reduces to the instance of the
crossing question at the column footprint: the forced column exists in every hypothetical tiling
(both corners break; a base-break consumes the single $c$), its descent is decided by the length
arithmetic except at the crossing, and a landing closes the family
(\texttt{A2BranchRow3.row\_three\_dies\_of\_span}, \texttt{span\_all\_a}).
\emph{(ii)} At $e\ge2$ the walk narrowing of Corollary~\ref{C-cor:basewalls} rests on
Theorem~\ref{C-thm:n1}, whose induction fails at the interior points $V_k$
(Remark~\ref{O-n1gap}); Remark~\ref{O-n1gapexact} shows the missing step \emph{is} the crossing
question at $V_k$, verbatim.
\emph{(iii)} Conjecture~\ref{C-conj:advance} is the crossing question by construction.
No fourth wall exists: the wedge-filler condition, once conjectured to block the regime
$e^4-4e^2f^2+2f^4\le0$, blocks nothing (the fill is forced at every member by the integrality of the
vertex system; \texttt{PrimeRegimes}).
\end{proposition}

\begin{proposition}[Nine tool classes cannot answer it]\label{prop:ninetools}\lab{PROVED}
Each of the following is closed with a recorded failure point: linear direction invariants and their
metric bounds (\texttt{LambdaFactor}); Euler and angle-sum counting, which coincide in a single
relation (\S\ref{sec:nogo}); edge-pairing parity; the non-abelian Conway--Lagarias tiling group,
which collapses to its rank-two abelianisation (Remark~\ref{OM-tilinggroup}); every criterion built
from similarity-invariant data of tile and target separately (\texttt{ScaleRigidity}, using only the
existence of the $m=2$ tiling); the $\gamma$-count on a stretch, whose two-endpoint hypothesis is
unsatisfiable under the monotone orientation word; the $b$-residue calculus, which constrains counts
exactly and locations not at all ($58$ of $60$ $b$-edges in the certificates have an endpoint that
is not a junction of the far side); De Bruijn--Hamel functionals, which require metric
incommensurability while every side of the family is an integer; and equilibrium stresses, whose
balance equations have slack at least two at every vertex type of the configuration, so that no
local forcing exists.
\end{proposition}

The common feature is exact: each tool computes a quantity invariant under relocating edges within
the tiling, and the crossing question asks \emph{where} an edge lies. The searches, which do decide
members, decide them precisely because a search tracks locations. Any future attack should be
measured against this requirement before effort is spent.

\section{Verification, claim by claim}\label{sec:verif}

The finite and numerical claims are checked by the scripts in the repository (exact arithmetic).
\texttt{verify\_shapes.py} reproduces the eleven-shape enumeration of Section~\ref{M-sub:shapes} and
the values $N_0$ of Section~\ref{M-sec:reduction}. \texttt{verify\_invariant.py} confirms
Lemma~\ref{M-lem:value} ($C_f(t)=\pm(c+a-b)$ over all orientations of $T$), Lemma~\ref{M-lem:cancel}
(both sides agree on explicit subdivision tilings, and~(the antisymmetry identity of \cite{Main}) cancels at a
non-edge-to-edge incidence), and Lemma~\ref{M-lem:nonint} (no counterexample among primitive
$120^\circ$-triples up to $c\approx 9\times10^6$). \texttt{verify\_spectrum.py} checks
Section~\ref{M-sec:full}: the two-positive-squares decomposition for all primes
$p\not\equiv3\pmod4$ below $10^5$; the scale structure $b\mid k^2\iff de\mid k$; the admissible
spectrum of Theorem~\ref{M-thm:admissible} (no prime is admissible; $33$ and $46$ are); and that
every instance of $e\mid(a+b-c)$ with $a,b<3000$ arises from the $j$-classification of
Proposition~\ref{M-prop:solv}, with none for $d=1$. As a positive control, both necessary conditions
of Theorem~\ref{M-thm:iso} hold on the genuine $2673$-tile isosceles tiling of
Herdt~\cite{BeesonIso}: for its tile $(5,3,7)$ and scale $k=27$, $N=k^2(a+2b)/b=2673$ and
$M=k(2b+a-2c)/(c+a-b)=-9$ are integers, as they must be on any true tiling.

The \emph{geometric} layer is also machine-checked, end to end and with no prose joints. That the
angles of the tiles meeting an interior point sum to $2\pi$ --- the fact a dissection argument uses
constantly and which has no counterpart in Mathlib --- is obtained without any theory of angular
measure, by a chain that is purely measure-theoretic:
\texttt{TangentCone.poly\_inter\_ball\_eq\_coneAt} (a polytope agrees with its tangent cone at a
point inside a small ball), \texttt{SectorArea.volume\_sector} (the unit sector of angle $\theta$
has area $\theta/2$, via polar coordinates), \texttt{Wedge.sector\_eq\_halfplanes} (that polar
sector is exactly the intersection of two half-planes with the unit disc),
\texttt{E2Join.volume\_halfplane\_wedge} and \texttt{E2Join.volume\_wedge} (the same area
statement transported to \texttt{EuclideanSpace} $\R$ \texttt{(Fin 2)}), and finally
\texttt{AngleSumAssembled.angle\_sum\_interior}. Every step is axiom-clean.

The entire \emph{arithmetic and combinatorial} layer of the proof is formalized and machine-checked
in Lean~4 with Mathlib (\texttt{lean/Erdos634/PrimeArithmetic.lean}, twenty-three theorems); every theorem depends only
on the standard axioms \texttt{propext}, \texttt{Classical.choice}, \texttt{Quot.sound}, with no
\texttt{sorry}:
\begin{itemize}
\item[] \texttt{shape\_enumeration}: the eleven-triple completeness of Section~\ref{M-sub:shapes} ---
realizable corner types with $\sum m_i=0$, $\sum k_i=3$ are exactly the eleven listed --- as pure
integer arithmetic;
\item[] \texttt{k\_not\_dvd\_sum\_sub}, \texttt{M\_not\_int}: Lemma~\ref{M-lem:nonint}
($(c-a-b)/\sqrt b\notin\Z$);
\item[] the full arithmetic of Theorem~\ref{M-thm:iso} in one master statement: no prime $N$
satisfies the $120^\circ$-relation, the area equation $Nb=k^2(a+2b)$, and the integrality
$(c+a-b)\mid k(2b+a-2c)$ simultaneously (\texttt{iso\_reduction\_identity},
\texttt{prime\_count\_forces\_scale}, \texttt{no\_prime\_isosceles\_count});
\item[] \texttt{add\_not\_prime}, \texttt{not\_prime\_of\_two\_le},
\texttt{F1\_count\_not\_prime}--\texttt{F4\_count\_not\_prime}: the scalene counts of
Proposition~\ref{M-prop:reduction} are composite ($a+b$ is never prime for a primitive
$120^\circ$-triple, and the $F_2,F_3,F_4$ counts factor);
\item[] \texttt{prime\_three\_mod\_four\_excluded}: a prime $p\equiv3\pmod4$ with $p>3$ is
neither a square, a sum of two squares, nor $2n^2$, $3n^2$ or $6n^2$ (the commensurable branch of
Section~\ref{M-sec:prelim}), via Fermat's two-squares theorem;
\item[] \texttt{prime\_sum\_two\_pos\_squares}: a prime $p\not\equiv3\pmod4$ is a sum of two
\emph{positive} squares (the achievability half of Theorem~\ref{M-thm:primefull});
\item[] \texttt{iso\_admissible}: Theorem~\ref{M-thm:admissible} --- given $b=de^2$ with $d$
squarefree, the area equation and $\Phi$-divisibility force $k=dew$, $N=dw^2(a+2b)$ and
$e\mid w(c-a-b)$;
\item[] \texttt{F1\_product}, \texttt{F1\_Ma}, \texttt{F1\_Mb}, \texttt{F1\_ratio},
\texttt{F1\_pin}, \texttt{disc\_not\_square}, \texttt{tile\_similar\_not\_prime}: the arithmetic of
Proposition~\ref{M-prop:product}, Corollary~\ref{M-cor:similar} and Proposition~\ref{M-prop:F1free}
(\texttt{lean/InvariantProduct.lean});
\item[] the finite arithmetic of the frontier sweeps, in eight theorems: the
$\pi/3$-equilateral square criterion fails at $14$ and $15$; the first tiling equation has no
admissible solution there; the two equilateral factor-pair uniqueness statements; the boundary
congruences that kill the isosceles candidates at $14$ and $22$; the $F_1$ invariant kill at
$21$; and the parity kill of $46$.
\end{itemize}
Mathlib still has no theory of planar dissections --- no \texttt{IsTiling}, no planar subdivision,
no Jordan curve theorem, no Euler formula, no polytopes --- so the geometric layer is built from the
convexity and measure-theoretic primitives that do exist. Within that constraint it is now
substantially machine-checked, and we state exactly how far, since the boundary has moved.

Formalized (\texttt{lean/Dissection.lean}). A dissection is defined as a target, $N$ tiles, the
covering equation, and pairwise disjointness of tile \emph{interiors} --- combinatorially, not
measure-theoretically, so that the measure-theoretic consequence is derived rather than assumed.
From it: \texttt{Dissection.aedisjoint} (disjoint interiors give a.e.\ disjointness, via
\texttt{Convex.addHaar\_frontier}; this is what lets a \emph{non-edge-to-edge} dissection still be
additive for area), the area identity \texttt{Dissection.volume\_target}, and
\texttt{vertex\_multiplicities}, which converts a real angle equation at a vertex into the integer
multiplicity equations the arithmetic consumes.

Also formalized: that each side of the target is a finite union of whole tile edges
(\texttt{hasEdgeChains\_edge}), by way of a support-face argument
(\texttt{SupportFace.mem\_convexHull\_max}: a point of a simplex attaining the maximum of a linear
functional is a convex combination of the vertices attaining it) together with a density step
(\texttt{SegmentDense.subset\_closure\_diff\_finite}).

Lemma~\ref{M-lem:cancel} was the one place where geometry remained, and it is now discharged.
\texttt{InvariantCore.cancellation\_core} derives the lemma from a single hypothesis: that the total
directed length of interior tile edges is invariant under reversal of direction. That hypothesis is
no longer assumed. \texttt{Dissection.g4\_final} proves it for every dissection, unconditionally.

The route avoids the notion of a maximal interior segment, for which Mathlib offers no support.
Instead of covering each maximal segment once from each side, one takes $S(d)$ to be the union of the
interior tile edges carrying unit direction $d$, and shows that $S(d)$ and $S(-d)$ agree away from the
finite set of tile vertices. The measure-theoretic ingredients are: a tile coincides with a half-plane
near a point in the relative interior of one of its edges (\texttt{Tri.inter\_ball\_eq\_halfplane}); a
half-plane through the centre halves a ball (\texttt{volume\_halfspace\_inter\_ball}, which Mathlib
does not have); a ball inside the target is partitioned by the tiles
(\texttt{Dissection.volume\_ball\_eq\_sum}); and the resulting count forces exactly two tiles at each
such point (\texttt{Dissection.two\_tiles\_at\_edge\_point}), the classification of a tile's local
contribution being exhaustive away from the tiles' vertices (\texttt{Tri.classify}), a finite set.

The orientation step carries the genuine geometric content, and runs as follows. The two tiles have
disjoint interiors, so the open half-planes they occupy near the shared point are disjoint
(\texttt{Dissection.cross\_disjoint\_of\_onEdge}, which uses the structural disjointness of interiors
rather than any measure); two nonzero linear functionals whose positive sets are disjoint are negative
multiples of one another (\texttt{proportional\_of\_disjoint\_pos}); hence the two edges span the
\emph{same} line with \emph{opposite} senses (\texttt{Dissection.leftDir\_antiparallel}); and
normalising to unit directions yields $S(d)\setminus F=S(-d)\setminus F$ with $F$ the finite vertex set
(\texttt{Dissection.dirSet\_horient}). Every statement named in this paragraph depends on
\texttt{propext}, \texttt{Classical.choice} and \texttt{Quot.sound}, and on nothing else.

What remains on written proofs, checked numerically but not formalized, is the direction grid of
Section~\ref{M-sec:prelim} and Lemma~\ref{M-lem:value}. Search verdicts (\texttt{EXHAUSTED}) are
declared computer assistance and are not formalized: doing so would mean verifying the search
engine. Finally, a formalized theorem certifies its own Lean statement; that each such statement
faithfully renders the paper result it is cited against is prose, and should be read as such.

\section*{Acknowledgements}

This work was carried out with substantial assistance from an artificial-intelligence system
(Anthropic's Claude), which proposed the signed-direction invariant, performed the symbolic and
numerical verification, found the elementary proof of Lemma~\ref{M-lem:nonint}, and produced the Lean
formalization, under the author's direction and review. The argument depends on the cited
classification of Laczkovich and Beeson for the branches other than the $2\pi/3$ tile, and on the
rationality theorem of Beeson and Zhang; the shape list coincides with Zhang's six sporadic
types~\cite{Zhang}. The present contribution is the invariant of Section~\ref{M-sec:invariant}, the
consequent exclusion of \emph{every} prime for the isosceles $2\pi/3$ branch (where Beeson had ruled
out only $N<36$ and left primality open), and the self-contained reduction of the scalene types. Its
arithmetic layer is machine-checked, and so now is the geometric layer, the interior angle sum
included (Section~\ref{M-sec:verification}). What is \emph{not} settled is the base-$\beta$ branch:
Theorem~\ref{M-thm:fullprime} is conditional on the complete-corner-wall hypothesis of the companion
note, and Remark~\ref{M-rem:conditional} states precisely what survives without it. The whole is
offered for independent refereeing.

\begin{thebibliography}{9}
\bibitem{Main} V.~Bonfioli, \emph{Erd\H{o}s Problem 634: the prime case}, 2026.
\bibitem{Companion} V.~Bonfioli, \emph{Erd\H{o}s Problem 634: the base-$\beta$ branch}, 2026.
\end{thebibliography}
\end{document}
